\documentclass[11pt,a4paper]{article}

\usepackage[utf8]{inputenc}
\usepackage[T1]{fontenc}
\usepackage{lmodern}   % scalable fonts (needed for microtype font expansion)
\usepackage{amsmath,amssymb,amsfonts}
\usepackage{mathtools}
\usepackage{bm}
\usepackage{geometry}
\usepackage{booktabs}
\usepackage{longtable}
\usepackage{array}
\usepackage{xcolor}
\usepackage{listings}
\usepackage{hyperref}
\usepackage{parskip}
\usepackage{enumitem}
\usepackage{microtype}

\geometry{margin=2.5cm}

\hypersetup{
  colorlinks=true,
  linkcolor=blue!60!black,
  urlcolor=blue!60!black,
  citecolor=blue!60!black
}

\lstset{
  basicstyle=\ttfamily\small,
  backgroundcolor=\color{gray!10},
  frame=single,
  framerule=0.4pt,
  breaklines=true,
  columns=fullflexible,
}

% Short-hands
\newcommand{\vp}{v_\perp}
\newcommand{\vpa}{v_\parallel}
\newcommand{\vpmin}{v_{\perp,\min}}
\newcommand{\vpmax}{v_{\perp,\max}}
\newcommand{\vpamin}{v_{\parallel,\min}}
\newcommand{\vpamax}{v_{\parallel,\max}}
\newcommand{\fvdf}{f(\vp,\vpa)}
\newcommand{\dt}{\Delta t}
\newcommand{\Lcal}{\mathcal{L}}
\newcommand{\namvar}[1]{\texttt{\textbf{#1}}}

% ============================================================
\title{\textbf{FP2D\_QLRF\_NL}\\[0.4em]
       \large User Manual\\[0.4em]
       \normalsize 2D Fokker--Planck solver for ion velocity distribution
       functions in cylindrical velocity space\\
       for one heated species or for all species of the plasma}
\author{F.~Louche\\LPP-ERM/KMS}
\date{Version 4.0 (multispecies) -- 24 November 2026}
% ============================================================

\begin{document}

\maketitle
\tableofcontents
\newpage

% ============================================================
\section{Introduction}
% ============================================================

\texttt{FP2D\_QLRF\_NL} solves the two-dimensional Fokker--Planck (FP) equation
for the velocity distribution function (VDF) $f(\vp,\vpa,t)$ of a minority ion
species in a magnetised plasma.  The two velocity coordinates are the speed
perpendicular to the magnetic field $\vp \ge 0$ and the speed parallel to the
field $\vpa \in (-\infty,+\infty)$.  Azimuthal symmetry about the magnetic
field direction is assumed throughout.

The code accounts for three physical operators:
\begin{enumerate}
  \item \textbf{Coulomb collisions} with one or more Maxwellian background
        species (electrons and bulk ions).
  \item \textbf{Quasi-linear radio-frequency (RF) heating} via the
        quasi-linear diffusion operator, evaluated using complex Bessel
        functions.
  \item \textbf{Self-collisions} (collisions of minority ions with
        themselves), with five levels of approximation selectable at runtime.
\end{enumerate}
An optional NBI \textbf{beam source} term and an associated particle
\textbf{loss term} are also included.

The code can compute either a \emph{steady-state} solution (direct solve) or
a \emph{time-dependent} evolution (implicit time integration).  Spatial
discretization uses a 7-point Fornberg finite-difference stencil.  The
resulting sparse linear systems are solved with Intel MKL PARDISO.

\paragraph{Version 4.0: multispecies.}  The main program is now \texttt{FP\_Coll\_2D\_ms}.  With its
defaults it is version 3.2, byte for byte.  With \texttt{i\_multi = -1} and
\texttt{i\_coupled = -1} it evolves \emph{every} species (heated species, electrons and bulk ions),
each on its own velocity grid, with the nonlinear collision operator of every pair built from
Rosenbluth potentials evaluated across grids by Legendre expansion.  Particle numbers, energy and
momentum are conserved to round-off by an implicit moment projection; an energy-confinement term
with the same $\tau_E$ on every species closes the balance when the plasma is driven by RF or by a
beam; the bulk grids can follow their species by regridding.  Section~\ref{sec:ms} describes this
mode, Sections~\ref{sec:nml_ms}, \ref{sec:out_ms} its namelist entries and outputs.  The rest of the
manual describes the single-species code, which is unchanged.


% ============================================================
\section{Physical Model}
% ============================================================

\subsection{The Fokker--Planck Equation}

The VDF $f = \fvdf$ evolves according to
\begin{equation}
\label{eq:FP}
  \frac{\partial f}{\partial t} =
    A\,f
  + B\,\frac{\partial f}{\partial \vp}
  + C\,\frac{\partial f}{\partial \vpa}
  + D\,\frac{\partial^2 f}{\partial \vp^2}
  + E\,\frac{\partial^2 f}{\partial \vp\,\partial \vpa}
  + F\,\frac{\partial^2 f}{\partial \vpa^2}
  + S(\vp,\vpa),
\end{equation}
where the six coefficient arrays $A$--$F$ (called \texttt{all00}, \texttt{all10},
\texttt{all01}, \texttt{all20}, \texttt{all11}, \texttt{all02} in the code) are
assembled as the sum of three independently computed contributions:
\begin{equation}
  X = X_{\mathrm{coll}} + X_{\mathrm{RF}} + X_{\mathrm{SC}},
  \qquad X \in \{A,B,C,D,E,F\},
\end{equation}
for Coulomb collisions, RF heating, and self-collisions, respectively.
The source term $S$ collects the NBI beam injection and the particle loss.

In steady-state mode ($\partial f/\partial t = 0$) the right-hand side of
\eqref{eq:FP} is set to zero and the resulting linear system is solved
directly.

\subsection{Coulomb Collision Operator}

The Coulomb collision operator is the linearized Fokker--Planck operator
describing the interaction of the minority ion species~$a$ (charge $Z_a$,
mass $m_a$, thermal velocity $v_a$) with a Maxwellian background species~$b$
(charge $Z_b$, mass $m_b$, thermal velocity $v_{Tb}$, density $n_b$).
The coefficient arrays $(\texttt{colin00}, \ldots, \texttt{colin02})$ for
each background species are computed by the subroutine \texttt{cblin} and
summed over all \texttt{nbulk} background species.

The collision frequency scale factor is
\begin{equation}
  \Gamma_b = \frac{n_b\,Z_a^2\,Z_b^2\,e^4\,\ln\Lambda_{ab}}
                  {4\pi\,\varepsilon_0^2\,m_a^2},
\end{equation}
where $\ln\Lambda_{ab}$ is the Coulomb logarithm evaluated at each time step
(see Section~\ref{sec:coulomb_log}).

\subsubsection{Coulomb Logarithm}\label{sec:coulomb_log}

Two distinct formulas are used:

\paragraph{Ion--electron:}
\begin{equation}
  \ln\Lambda_{ae} = 24 - \ln\!\left(\frac{n_e^{1/2}}{T_e}\right),
\end{equation}
with $T_e$ in eV and $n_e$ in m$^{-3}$.

\paragraph{Ion--ion:}
\begin{equation}
  \ln\Lambda_{ab} = 23 - \ln\!\left[
    Z_a Z_b\,\frac{m_a+m_b}{m_a\,T_b + m_b\,T_a}\,
    \left(\frac{n_a\,Z_a^2}{T_a}+\frac{n_b\,Z_b^2}{T_b}\right)^{1/2}
  \right].
\end{equation}
In time-dependent runs the Coulomb logarithm is recomputed at each time step
using the current effective temperature of the minority ion species.  (Versions before 4.0 used
$Z_a^2Z_b^2$ in the mass factor; the two agree for $Z=1$.)  The multispecies mode adds the
electron--electron logarithm (Section~\ref{sec:ms}).

\subsection{Quasi-Linear RF Operator}

When radio-frequency heating is active (\texttt{irf} $= -1$), the
quasi-linear diffusion operator is added to the coefficient arrays.  The RF
operator describes the resonant power exchange between the minority ions and
an electromagnetic wave characterised by:
\begin{itemize}
  \item Wave frequency $\omega = 2\pi f$ (\texttt{frek} in Hz).
  \item Parallel wave vector $k_\parallel$ (\texttt{kpar} in rad/m).
  \item Perpendicular wave vector $k_\perp$ (complex, \texttt{kperp}
        in rad/m).
  \item Left- and right-hand polarisation components $E_+$, $E_-$
        (\texttt{eplus}, \texttt{emin}).
  \item Static magnetic field $B_0$ (\texttt{b0} in T).
  \item Harmonic number $n$ (\texttt{nharm}).
\end{itemize}

The quasi-linear diffusion coefficients are computed in the subroutine
\texttt{qlrfterm} using complex Bessel functions $J_n(k_\perp \vp / \omega_c)$
and their derivatives, evaluated at each grid point.  The resonance condition
is $\omega - n\,\omega_c - k_\parallel\,\vpa = 0$, broadened by a Gaussian
of width \texttt{delta\_RF} in $\vpa$.

\subsection{Self-Collision Operator}

Four options are available for the treatment of self-collisions, selected by
the flag \texttt{isc}:

\paragraph{\texttt{isc = 0} -- No self-collisions.}
The self-collision coefficients are zero.  The FP equation is linear.

\paragraph{\texttt{isc = 1} -- Maxwellian background at fixed temperature.}
The self-collision operator is computed using the Rosenbluth potentials for a
Maxwellian distribution at the Stix effective temperature $T_\mathrm{eff}$
(computed once at the start of the simulation and kept constant).  This adds
a fixed linear operator to the coefficient arrays.

\paragraph{\texttt{isc = 2} -- Maxwellian background at varying effective temperature.}
Same as \texttt{isc = 1}, but the background temperature and Coulomb
logarithm are updated at each time step to track the evolving effective
temperature $T_\mathrm{eff}(t) = (2T_\perp + T_\parallel)/3$.  Requires a
time-dependent run.

\paragraph{\texttt{isc = 3} -- Maxwellian background at the density-characteristic
temperature.}
Same as \texttt{isc = 2}, but the background temperature is the
\emph{density-characteristic} (cold-bulk) temperature $T_n(t)$ instead of the
energy-weighted $T_\mathrm{eff}(t)$.  $T_n$ is obtained from a
phase-space-weighted least-squares fit of $\ln f$ versus $v^2$ restricted to the
dense thermal core --- only cells with $f > \texttt{core\_frac}\cdot\max f$ enter
the fit --- which returns the temperature exactly for a Maxwellian.  Being tied to
the cold bulk rather than to the supra-thermal tail, this background better
represents the population that dominates the collisional drag.  The core fraction
\texttt{core\_frac} (default $3.8\times10^{-3}$) is tunable; the appropriate value
depends on the shape of the distribution, so it may need re-calibrating against the
fully nonlinear operator for a new scenario.  Requires a time-dependent run.

\paragraph{\texttt{isc = -1} -- Fully nonlinear self-collisions.}
The complete nonlinear Landau/Fokker--Planck self-collision operator
\begin{equation}
  \left(\frac{\partial f}{\partial t}\right)_{\mathrm{SC}}
  = \Gamma_{aa}\,\frac{\partial}{\partial v_i}
    \left[ H_{ij}[f]\,\frac{\partial f}{\partial v_j}
         - G_i[f]\,f \right],
\end{equation}
is used, where $H_{ij}$ and $G_i$ are the Rosenbluth potentials built from
the \emph{current} distribution $f(t)$.  The Rosenbluth potentials are
evaluated via a precomputed phi-distance kernel \texttt{sum\_phi} using
Gauss--Legendre quadrature (subroutine \texttt{distance\_v\_gauss\_legendre}).
This kernel depends only on the velocity grid and is computed once and saved
to disk (\texttt{sum\_phi.dat}); it can be reloaded for subsequent runs on the
same grid.  The operator is rebuilt at every time step, making the problem
nonlinear.  Requires a time-dependent run.

\subsection{Beam Source and Particle Losses}\label{sec:beam}

When \texttt{isource} $= -1$, a neutral-beam injection (NBI) source is added:
\begin{equation}
  S(\vp,\vpa) = \frac{n_0}{\tau_s}\,
    \frac{\delta_G(\vp - \vp^\mathrm{beam})\,\delta_G(\vpa - \vpa^\mathrm{beam})}
         {\mathcal{N}},
\end{equation}
where $\delta_G(x) = e^{-(x/\sigma)^2}/(\sqrt{\pi}\,\sigma)$ is a Gaussian
approximation to the Dirac delta, $\vp^\mathrm{beam} = v_\mathrm{beam}\sin\alpha$,
$\vpa^\mathrm{beam} = v_\mathrm{beam}\cos\alpha$ with $\alpha$ the injection
angle, and $\mathcal{N}$ is the normalization factor
\begin{equation}
  \mathcal{N} = \sqrt{\pi}\,\sigma_\perp\,e^{-(\vp^\mathrm{beam}/\sigma_\perp)^2}
               + \pi\,\vp^\mathrm{beam}\,\mathrm{erfc}(-\vp^\mathrm{beam}/\sigma_\perp),
\end{equation}
which reduces to $2\pi\,\vp^\mathrm{beam}$ for $\vp^\mathrm{beam} \gg \sigma_\perp$
and to $\sqrt{\pi}\,\sigma_\perp$ for a purely parallel beam ($\vp^\mathrm{beam}=0$).

The loss term is $-f/\tau_s$ (isotropic particle losses with slowing-down time $\tau_s$).


% ============================================================
\section{Numerical Methods}
% ============================================================

\subsection{Velocity-Space Grid}

The computational domain is $[\vpmin,\,\vpmax] \times [\vpamin,\,\vpamax]$,
discretized on a $N_\perp \times N_\parallel$ grid with 1-based indices
$(i,j)$ where $i \in [1,N_\perp]$ and $j \in [1,N_\parallel]$.

Three grid spacing modes are available, controlled by \texttt{ising}:

\begin{description}
  \item[\texttt{ising = 0}] Uniform spacing in both directions.
        $\Delta\vp = (\vpmax-\vpmin)/(N_\perp-1)$,
        $\Delta\vpa = (\vpamax-\vpamin)/(N_\parallel-1)$.

  \item[\texttt{ising = -1}] Two-domain grid in $\vp$.  A finer mesh of
        \texttt{nsing} points is used for $\vp \le \texttt{vbound}$, and a
        coarser uniform mesh for $\vp > \texttt{vbound}$.  This provides
        higher resolution near the magnetic axis where gradients are typically
        steeper.

  \item[\texttt{ising = +1}] Power-law spacing in $\vp$,
        ${\vp}_i \propto \bigl[(i-\tfrac12)/(\texttt{nperp}-1)\bigr]^{p}$ with
        $p = \texttt{p\_grid}$ (default 2, quadratic), providing higher
        resolution near the axis $\vp = 0$.  The \texttt{fd\_stencil\_2d}
        routine uses the Fornberg algorithm (\texttt{fornberg\_weights}) to
        compute exact finite-difference weights for the local non-uniform
        spacing, so this mode is fully supported.
\end{description}

The cylindrical Jacobian is $J = 2\pi\,\vp$, so that integrals over velocity
space are $\int\!\int f\,J\,\mathrm{d}\vp\,\mathrm{d}\vpa$.

The global 1D degree-of-freedom (DOF) index is
\begin{equation}
  \mathrm{ix}(i,j) = (i-1)\,N_\parallel + j,
  \quad i \in [1,N_\perp],\; j \in [1,N_\parallel],
\end{equation}
giving $N_\mathrm{dof} = N_\perp \times N_\parallel$ total DOFs.

\subsection{Finite-Difference Discretization}

Spatial derivatives are discretized using the \textbf{7-point Fornberg
stencil} (subroutine \texttt{fd\_stencil\_2d}), which is a fourth-order
accurate scheme.  The stencil uses 4 points in each direction: at the current
node and its 3 neighbours.  For a non-uniform $\vp$ grid, the Fornberg
algorithm computes the exact derivative weights for the local grid spacing via
subroutine \texttt{fornberg\_weights}.

The stencil returns at most 49 entries per row (7~points in $\vp$ by 7 in
$\vpa$), although boundary rows have fewer entries.

\subsubsection{Boundary Conditions}

\begin{description}
  \item[Axis ($i=1$, $\vp=\vpmin$):] Neumann condition $\partial f/\partial\vp = 0$,
        enforcing axial symmetry.

  \item[Outer boundary ($i=N_\perp$):] Dirichlet condition $f = 0$.

  \item[Parallel boundaries ($j=1$ and $j=N_\parallel$):] Dirichlet condition $f = 0$.
\end{description}

\subsection{Advection Stabilization at High $\vp$ (\texttt{i\_upwind})}
\label{sec:upwind}

The perpendicular drag coefficient $B$ (collisional friction) grows with $\vp$
while the perpendicular diffusion $D$ falls off on the supra-thermal tail, so the
cell-P\'eclet number
\begin{equation}
  \mathrm{Pe}_\perp = \frac{|B|\,\Delta\vp}{D}
\end{equation}
increases steeply toward the outer boundary and can reach $\mathrm{Pe}_\perp \sim
40$ at $\vp = \vpmax$ for high-energy (e.g.\ second-harmonic ICRF) scenarios.
Beyond $\mathrm{Pe}_\perp \gtrsim 2$ a central discretization of the advection
term is no longer dissipative; at the outer Dirichlet wall, where the Fornberg
stencil is one-sided, the effective perpendicular diffusion of the discrete
operator can even become \emph{negative}.  A modified-equation analysis of the
Dirichlet-truncated wall row gives, in the uniform-spacing limit,
\begin{equation}
  D_{\mathrm{eff}} \;=\; \frac{223}{360}\,D \;-\; \frac{1}{12}\,|B|\,\Delta\vp ,
\end{equation}
i.e.\ the discrete diffusion turns negative exactly when the wall-cell P\'eclet
number exceeds $223/30 \approx 7.4$ (the deleted wall node biases the drag stencil
downwind; high order shrinks the classic two-point downwind coefficient from
$\tfrac12$ to $\tfrac{1}{12}$ but cannot change its sign).  The assembled operator $\Lcal$ then
acquires a small set of eigenvalues with positive real part, localized in a thin
layer against the $\vp = \vpmax$ boundary.  Under Crank--Nicolson
(\texttt{icn = -1}) these modes are amplified and the run eventually blows up,
whereas the fully-implicit scheme damps them.  The instability is governed by
$\vpmax$ (the wall velocity) and is essentially independent of the RF term and of
the grid resolution.

Setting \texttt{i\_upwind = -1} cures this at the discretization level: in every
cell with $\mathrm{Pe}_\perp > 2$ the drag term $B\,\partial f/\partial\vp$ is
evaluated with a first-order \emph{upwind} one-sided difference (forward for
$B>0$, backward for $B<0$, i.e.\ the side that makes the semi-discrete eigenvalue
dissipative), while the diffusion term keeps its central Fornberg weights.  In the
modified equation the upwind drag moment $+|B|\,\Delta\vp/2$ replaces
$-|B|\,\Delta\vp/12$, so $D_{\mathrm{eff}} > 0$ at \emph{any} P\'eclet number; the
$\mathrm{Pe}_\perp > 2$ switching threshold is therefore conservative with respect
to the $\approx 7.4$ instability threshold above.  The
sparsity pattern is unchanged (the diffusion stencil already fills every
same-column entry), so the phased-PARDISO machinery is untouched.  This removes
all positive-real-part modes with a bulk-physics change of at most a few percent,
because the upwinding acts only on the near-empty tail.  The default
\texttt{i\_upwind = 0} keeps the pure central scheme; running fully implicit
(\texttt{icn = 0}) is an alternative stabilization.  Note that first-order
upwinding adds some numerical diffusion on the tail; a flux-limited or
exponentially-fitted variant would reduce this if higher tail accuracy is needed.

\subsection{Sparse Matrix Storage}

All matrices are stored in \textbf{1-based Compressed Sparse Row (CSR)} format.
For a matrix of size $N_\mathrm{dof}$:
\begin{itemize}
  \item \texttt{ia(1:$N_\mathrm{dof}$+1)}: row pointer array.
  \item \texttt{ja(1:nnz)}: column indices.
  \item \texttt{aa(1:nnz)}: non-zero values.
\end{itemize}
Entries within each row are sorted by column index before passing to PARDISO.
The maximum number of non-zeros per row is 49 (7$\times$7 stencil), so
$\texttt{nnz\_max} = 49\,N_\mathrm{dof}$.

\subsection{Direct Solver: Intel MKL PARDISO}

All linear systems are solved with \textbf{Intel MKL PARDISO} using matrix
type \texttt{mtype = 11} (real non-symmetric general).  The module
\texttt{pardiso\_solver} exposes four phased entry points that exploit
PARDISO's factorization phases:

\begin{center}
\begin{tabular}{ll}
\toprule
Entry point & PARDISO phases \\
\midrule
\texttt{pardiso\_solve\_steady}  & 11 + 22 + 33 (full solve, then release) \\
\texttt{pardiso\_solve\_init}    & 11 + 22 (symbolic + numeric factorization) \\
\texttt{pardiso\_solve\_step}    & 33 only (back-substitution) \\
\texttt{pardiso\_solve\_finalize}& release internal memory \\
\bottomrule
\end{tabular}
\end{center}

This phased approach is critical for performance in the time-dependent solver:
when the LHS matrix is constant (linear problem), the factorization is done
once and only the back-substitution is repeated at each time step.


% ============================================================
\section{Algorithms}
% ============================================================

\subsection{Steady-State Solver}

When \texttt{ntimes(1) = 0}, the code solves the steady-state FP equation
\begin{equation}
  \Lcal[f] = -S,
\end{equation}
where $\Lcal$ is the FP differential operator assembled from the coefficient
arrays \texttt{all00}, \ldots, \texttt{all02}.  The two-pass assembly (first
count non-zeros, then fill values) builds the CSR matrix, which is then
passed to \texttt{pardiso\_solve\_steady} for a direct solve (subroutine
\texttt{FP\_steady\_state} in module \texttt{mod\_linear}).

\subsection{Time-Dependent Solver -- Linear Problem (\texttt{timefp\_7pt})}

Used when \texttt{ntimes(1) > 0} and $\texttt{isc} \ne -1$.

\subsubsection{Time Discretization}

A $\theta$-method is used:
\begin{equation}
  \bigl(I - \theta\,\dt\,\Lcal\bigr)\,f^{n+1}
  = \bigl(I + (1-\theta)\,\dt\,\Lcal\bigr)\,f^n + \dt\,S.
\end{equation}

\begin{center}
\begin{tabular}{lll}
\toprule
\texttt{icn} & $\theta$ & Scheme \\
\midrule
$-1$  & $0.5$  & Crank--Nicolson \\
$+1$  & $0.75$ & Intermediate \\
else  & $1.0$  & Fully implicit \\
\bottomrule
\end{tabular}
\end{center}

\subsubsection{Algorithm}

\begin{enumerate}
  \item \textbf{Assemble} the FP operator $\Lcal$ in CSR format (two-pass loop).
  \item \textbf{Build LHS}: $M_\mathrm{lhs} = I - \theta\,\dt\,\Lcal$.  The sparsity
        pattern is identical to $\Lcal$ since the diagonal already exists.
  \item \textbf{Factorize} $M_\mathrm{lhs}$ once (PARDISO phases 11+22).
  \item \textbf{Time loop} for $n = 1, \ldots, N_\mathrm{times}$:
        \begin{enumerate}
          \item Recompute the Coulomb logarithm and, for
                \texttt{isc} $\in \{1,2,3\}$, the Maxwellian self-collision
                background (at the fixed Stix, $T_\mathrm{eff}(t)$, or $T_n(t)$
                temperature, respectively) from the current distribution, then
                rebuild the FP operator (allowing a time-varying Coulomb logarithm
                and self-collision term).
          \item Rebuild $M_\mathrm{lhs}$ with the updated operator.
          \item Compute the RHS:
                $b = \bigl(I + (1-\theta)\,\dt\,\Lcal\bigr)\,f^n + \dt\,S$
                via a hand-coded sparse matrix--vector product.
          \item Enforce zero RHS on boundary DOFs.
          \item Solve $M_\mathrm{lhs}\,f^{n+1} = b$ (PARDISO phase 33 only).
          \item Compute diagnostics (density, energy, power, momentum).
        \end{enumerate}
  \item \textbf{Finalize} PARDISO.
  \item If \texttt{isource = 0}: rescale $f \leftarrow f \cdot n_0/\int f\,\mathrm{d}^3v$
        to conserve total particle number.
\end{enumerate}

\subsection{Time-Dependent Solver -- Nonlinear Self-Collisions (\texttt{timefp\_7pt\_nl})}

Used when \texttt{ntimes(1) > 0} and \texttt{isc = -1}.  The LHS matrix now
changes at every time step because the self-collision operator depends on
$f^n$.

\subsubsection{Phi-Distance Kernel Precomputation}

Before entering the time loop, the phi-distance kernel
$\Phi_{ij} = \phi(|\bm{v}_i - \bm{v}_j|)$ (Gauss--Legendre quadrature,
$O(N_\mathrm{dof}^2)$) is cached on disk in a per-case file
\texttt{sum\_phi-CASENAME.dat} (or \texttt{sum\_phi.dat} when \texttt{casename}
is empty). Since the kernel is grid-dependent, the file header stores a grid
signature (\texttt{nperp}, \texttt{npar}, \texttt{vperp\_min/max},
\texttt{vpar\_min/max}). Caching is fully self-managing --- no user flag is
required:

\begin{center}
\begin{tabular}{ll}
\toprule
Cache file state & Action \\
\midrule
Missing                                  & Compute kernel and save \\
Present, signature matches               & Load kernel (skip the quadrature) \\
Present, signature differs or truncated  & Recompute and overwrite \\
\bottomrule
\end{tabular}
\end{center}

Because the signature is verified automatically, a stale kernel from a
different velocity grid can never be silently reused.

\subsubsection{Time-Stepping Scheme}
\label{sec:nl_theta}

The $\theta$ parameter has an additional option for the NL solver:

\begin{center}
\begin{tabular}{lll}
\toprule
\texttt{icn} & $\theta$ & Comment \\
\midrule
$-1$ & $0.5$ & Crank--Nicolson (may be unstable at high $\vp$; see Sec.~\ref{sec:upwind}) \\
$+1$ & $0.75$ & Intermediate \\
else & $1.0$ & Fully implicit (recommended) \\
\bottomrule
\end{tabular}
\end{center}

\subsubsection{Algorithm}

\begin{enumerate}
  \item \textbf{Build sparsity pattern} of $\Lcal$ using the linear
        coefficients only (pattern is fixed; NL SC uses the same stencil).
  \item \textbf{Factorize symbolically} (PARDISO phase 11 only).
  \item \textbf{Time loop} for $n = 1, \ldots, N_\mathrm{times}$:
        \begin{enumerate}
          \item \textbf{Build NL self-collision term}: call
                \texttt{main\_nlterm}$(f^n)$ to compute the self-collision
                coefficient arrays \texttt{sc**} from the current distribution
                via the phi-distance kernel.
          \item \textbf{Update Coulomb logarithm} and rebuild the linear FP
                coefficients.
          \item \textbf{Assemble total operator}: $\Lcal = \Lcal_\mathrm{lin} + \Lcal_\mathrm{SC}$.
          \item \textbf{Rebuild LHS}: $M_\mathrm{lhs} = I - \theta\,\dt\,\Lcal$.
          \item \textbf{Compute RHS}:
                $b = (I+(1-\theta)\,\dt\,\Lcal)\,f^n + \dt\,S$.
          \item \textbf{Solve} (PARDISO phases 22+33 every step since $M_\mathrm{lhs}$ changes).
          \item Compute diagnostics.
        \end{enumerate}
  \item \textbf{Finalize} PARDISO.
\end{enumerate}

\textit{Note on start-up:} when starting from $f = 0$ with a beam source
(i.e., \texttt{istart = 0}), the NL self-collision term is skipped for the
first few steps until the beam density becomes significant enough for the
Rosenbluth potentials to be numerically meaningful.

\subsection{Multi-Phase Time-Stepping}

The namelist parameters \texttt{ntimes} and \texttt{timestep} are each
3-element arrays, enabling up to three sequential phases with different time
steps:
\begin{itemize}
  \item Phase~$k$ runs \texttt{ntimes(k)} steps of duration \texttt{timestep(k)} seconds.
  \item Phases with \texttt{ntimes(k) = 0} are skipped.
  \item Phases 2 and 3 append to the output files (restart in time).
\end{itemize}
Old single-phase input files (\texttt{ntimes = N, timestep = dt}) remain
backward-compatible: the scalar value fills element~1 and elements 2--3
default to zero.

\subsection{Initial Condition for Time-Dependent Runs}

The initial distribution $f^0$ is selected by \texttt{istart}:

\begin{center}
\begin{tabular}{cl}
\toprule
\texttt{istart} & Initial condition \\
\midrule
0 & $f^0 = 0$ everywhere. Requires \texttt{isource = -1}. \\
1 & Stix Maxwellian solution (default). \\
2 & Steady-state solution of the linear solver without self-collisions. \\
3 & Steady-state solution of the linear solver with Maxwellian SC background. \\
\bottomrule
\end{tabular}
\end{center}

\subsection{Steady-State Convergence Detection}
\label{sec:ssconv}

When \texttt{i\_ss\_check = -1} the run stops early once the solution stops
changing.  A \emph{single} test decides this, implemented in
\texttt{mod\_conv\_diag}.  It is evaluated once every $W =
\texttt{n\_ss\_window}$ steps -- not at every step -- and measures rates of
change over that window, exactly as the former scalar criterion did, but on the
distribution function itself rather than on four integrated scalars.  At each
evaluation the current state is compared with the state stored at the previous
evaluation, so only \emph{one} snapshot of $f$ is retained regardless of $W$,
and the full-grid norms are formed once per window rather than once per step.

\paragraph{Phase-space weighting.}  $f$ lives in cylindrical velocity space, so
every norm and moment is integrated against $\mathrm{d}^3v = 2\pi\vp\,
\mathrm{d}\vp\,\mathrm{d}\vpa$.  Composite-trapezoidal weights on the (possibly
non-uniform) grids are folded with the Jacobian into $W_\perp(i) =
2\pi v_{\perp,i}\, w_\perp(i)$ and $W_\parallel(j) = w_\parallel(j)$.  Ignoring
the Jacobian would over-weight the near-axis region, which holds almost no
particles.  As a check that both the Jacobian and the non-uniform weights are
right, the code verifies the exact identity
$\sum_i W_\perp(i) = \pi(v_{\perp,N}^2 - v_{\perp,1}^2)$ at startup.

\paragraph{Criteria.}  Writing $f^\mathrm{now}$ for the state at the current
evaluation and $f^\mathrm{prev}$ for the one at the previous evaluation, and
with $\mathrm{d}t_\mathrm{win} = t_\mathrm{now} - t_\mathrm{prev}$ the physical
duration of the window, the test computes
\begin{align}
  \varepsilon      &= \frac{\Vert f^\mathrm{now} - f^\mathrm{prev}\Vert_2}
                           {\mathrm{d}t_\mathrm{win}\,\Vert f^\mathrm{now}\Vert_2}, &
  \varepsilon_\mathrm{tail} &= \frac{\Vert f^\mathrm{now} - f^\mathrm{prev}\Vert_T}
                           {\mathrm{d}t_\mathrm{win}\,\Vert f^\mathrm{now}\Vert_T},
  \label{eq:epsconv}
\end{align}
where $\Vert\cdot\Vert_2$ uses the weights above and $\Vert\cdot\Vert_T$ adds the
tail weight $T = (1 + v^2/v_\mathrm{th}^2)^{k}$ ($k=2$).  The Jacobian-weighted
norm is dominated by the bulk, so $\varepsilon_\mathrm{tail}$ is reported
\emph{separately}: the interesting failure mode is $\varepsilon$ small while
$\varepsilon_\mathrm{tail}$ is not, i.e.\ a converged bulk over a still-evolving
tail.  The relative drift rates of the density, parallel flow and energy are
computed over the same window; norm convergence is necessary but not sufficient,
and a moment still drifting while $\varepsilon$ is small points to a flux or
boundary-condition error at $\vp=0$ or at the domain edges rather than to a
convergence failure.

With the default \texttt{i\_conv\_shape = 0} the run stops when \emph{all five}
conditions hold simultaneously:
$\varepsilon/\nu_\mathrm{ref} < \texttt{ss\_tol\_eps}$,
$\varepsilon_\mathrm{tail}/\nu_\mathrm{ref} < \texttt{ss\_tol\_tail}$, and each
of the three moment drifts $< \texttt{ss\_tol\_moment}$.

\paragraph{Amplitude vs shape (\texttt{i\_conv\_shape}).}  The rates above measure
the change of $f$ \emph{including its amplitude}, which is the right thing when the
run really does reach a steady state.  In a \emph{sourceless} run it is not: particles
absorbed at the Dirichlet boundaries are never replaced, so once the VDF shape has
settled $f$ keeps decaying as $\phi(v)\,e^{-t/\tau}$ with the shape frozen.  Then
\begin{equation}
  \varepsilon \;\longrightarrow\; \frac{e^{\Delta t_\mathrm{win}/\tau}-1}
                                       {\Delta t_\mathrm{win}} \;\approx\; \frac{1}{\tau},
\end{equation}
a floor that no tolerance below $1/\tau$ can ever clear, and $\varepsilon$,
$\varepsilon_\mathrm{tail}$ and the density and energy drifts all collapse onto that
same value (the tell-tale signature of a pure rescaling).

Setting \texttt{i\_conv\_shape = -1} selects the \emph{shape} criterion: each field is
normalised to unit norm before differencing,
\begin{equation}
  \varepsilon_\mathrm{shape} =
    \frac{1}{\Delta t_\mathrm{win}}
    \left\Vert \frac{f^\mathrm{now}}{\Vert f^\mathrm{now}\Vert}
             - \frac{f^\mathrm{prev}}{\Vert f^\mathrm{prev}\Vert} \right\Vert ,
  \label{eq:epsshape}
\end{equation}
so a uniform rescaling $f^\mathrm{now} = c\,f^\mathrm{prev}$ gives exactly zero.  The
moment tests then switch to the \emph{intensive} moments --- the parallel flow
$u_\parallel$ and the mean energy per particle $E/n$, both unchanged by a uniform
drain --- and the density drift is dropped, being pure amplitude.  All four rates are
always computed and reported; \texttt{i\_conv\_shape} only selects which family the
tolerances are applied to, and the log line is tagged \texttt{[amplitude]} or
\texttt{[shape]} accordingly.

Verified 2026-07-30 on a 60$\times$41 \texttt{isc = 2} RF case: the amplitude rates
flattened at $\varepsilon \approx 2\times10^{-3}$ (the drain rate, with the density
drift pinned at $1.73\times10^{-3}$) and the run never stopped, while
$\varepsilon_\mathrm{shape}$ kept falling exponentially and reached
$8.7\times10^{-5}$ at $t = 12.5$~s, where the shape criterion stopped the run.

\paragraph{Units.}  $\varepsilon$ has dimensions of inverse time, so it is only
meaningful against a physical rate; it is divided by a reference rate
$\nu_\mathrm{ref}$ (the code passes $1/\tau_{ie}$) before comparison, and the
tolerances are therefore dimensionless fractions of that rate.  Because the
window duration is taken from the simulation clock rather than a step count, the
tolerances are \emph{independent of the time step} and remain correct across
phase boundaries where \texttt{timestep} changes.

\paragraph{Cost.}  One snapshot of $f$ is retained
($n_\perp n_\parallel \times 8$~bytes, independent of $W$; the figure is printed
at startup), and one full-grid pass is performed per window, i.e.\ $1/W$ of the
cost of testing at every step.

Diagnostics are printed and written once per window -- to
\texttt{conv\_eps\_vs\_time.txt} (time, $\varepsilon$,
$\varepsilon_\mathrm{tail}$, $\varepsilon_\mathrm{shape}$,
$\varepsilon_\mathrm{tail,shape}$) and \texttt{conv\_diag\_vs\_time.csv}
(those four rates, their ratios to $\nu_\mathrm{ref}$, the moments and all the
drift rates).  The first evaluation only lays down the reference state, so the
first verdict comes at step $2W$.

\subsection{Diagnostics and Output}

At each time step the following quantities are written to disk:
\begin{itemize}
  \item Density $n(t) = \int f\,J\,\mathrm{d}^3v$.
  \item Total, perpendicular, and parallel kinetic energies.
  \item Effective temperature $T_\mathrm{eff} = (2T_\perp + T_\parallel)/3$,
        where $T_\perp = \tfrac{1}{2}m\langle \vp^2 \rangle$ and
        $T_\parallel = m\langle \vpa^2 \rangle$.
  \item Density-characteristic temperature $T_n$ (the cold-bulk temperature from
        the $\ln f$-vs-$v^2$ log-slope fit; it is the background used by
        \texttt{isc = 3} but is written as a diagnostic for \emph{all} runs).
  \item Near-axis and whole-grid minima of $f$ (a negative-$f$ monitor: small
        negative undershoots near the axis can corrupt low-velocity moments such
        as $T_n$).
  \item Perpendicular anisotropy factor $= 100\,T_\perp/(T_\perp + T_\parallel)$
        (\%), ranging from 0\% (all-parallel) to 50\% (Maxwellian) to 100\%
        (all-perpendicular).
  \item Power densities (collisional by species, RF, beam source, losses,
        self-collision total), and their perpendicular/parallel split for
        self-collisions.
  \item Momentum transfer rates (perpendicular and parallel) for each
        operator.
  \item Coulomb logarithm(s) vs.\ time.
\end{itemize}

At the end of a run, balance diagnostics (power, density, and optionally
momentum) and a check of the size of the velocity domain are printed to
standard output (Section~\ref{sec:endofrun}).

% ============================================================
\section{Multispecies Mode (Version 4.0)}
\label{sec:ms}
% ============================================================

Version 4.0 can evolve \emph{every} species of the plasma, not only the heated one. The electrons
and each bulk ion species then have their own velocity grid and distribution, collide with each
other and with the heated species through the full nonlinear Landau operator, and can lose energy
through a confinement term. This section describes the model, the algorithm and their use; the
namelist entries are listed in Section~\ref{sec:nml_ms} and the outputs in
Section~\ref{sec:out_ms}. The method and its verification are documented in detail in the
multispecies article (\emph{Multispecies extension of NLFP2D with species-specific grids}, living
document in the development repository).

\subsection{Three Levels of Use}

\begin{longtable}{>{\raggedright}p{4.2cm} p{10.4cm}}
\toprule
Namelist & What is computed\\\midrule
\endhead
\bottomrule
\endlastfoot
\texttt{i\_multi = 0} (default) & Version 3.2, unchanged: the heated species against fixed Maxwellian
  backgrounds.  With $\tau_E>0$ and a source, $L_E$ (below) acts on the heated species. \\[3pt]
\texttt{i\_multi = -1} & As above, but a container, grid and initial Maxwellian are also built for
  every background species and written to \texttt{finit-*.txt}.  The backgrounds stay frozen; this
  level serves the checks of Section~\ref{sec:ms_checks}. \\[3pt]
\texttt{i\_multi = -1}, \texttt{i\_coupled = -1} & \textbf{Coupled multispecies run}: every species
  evolves on its own grid (Section~\ref{sec:ms_loop}).  Requires \texttt{isc = -1}, \texttt{iold = 0}
  and a time-dependent run. \\
\end{longtable}

\subsection{Coupled Equations}

The species are the heated species $a$ (\texttt{aa}, \texttt{za}, density
$\texttt{xpart}\cdot n_e$), the electrons and the \texttt{nbulk}$-1$ bulk ion species (\texttt{ab},
\texttt{zb}, \texttt{xb}).  Each gyrotropic distribution $f_s(\vp,\vpa,t)$ obeys
\begin{align}
\frac{\partial f_a}{\partial t}&=\sum_jC(f_a,f_j)+Q_\mathrm{RF}(f_a)+S_\mathrm{NBI}
  -\frac{f_a}{\tau_s}\Big|_\mathrm{NBI}+L_E(f_a),\\
\frac{\partial f_b}{\partial t}&=\sum_jC(f_b,f_j)+L_E(f_b),\qquad b\in\{e,i_1,i_2,\dots\},
\end{align}
where $j$ runs over all species, self included.  Only the heated species interacts with the RF
field or is injected by the beam.  The densities of the bulk species are fixed by the namelist
(quasi-neutrality is the user's responsibility).

\paragraph{Collision operator of a pair.}
In a pair $(a\leftarrow b)$, $a$ is the \emph{target} and $b$ the \emph{background}, the species whose
Rosenbluth potentials act on $a$.  With $\psi_b=-g_b/8\pi$, $\phi_b=\nabla^2\psi_b=-h_b/4\pi$ and
$\mu=m_a/m_b$,
\begin{equation}
C(f_a,f_b)=\kappa_{ab}\Big[\psi_{b,ij}\,f_{a,ij}+(1-\mu)\,\phi_{b,j}\,f_{a,j}-\mu\,f_b\,f_a\Big],\qquad
\kappa_{ab}=-4\pi\gamma_0\ln\Lambda_{ab}\Big(\frac{Z_a}{A_a}\Big)^2Z_b^2,
\end{equation}
with $\gamma_0=0.2390775$; for $\mu=1$ this is the v3.2 self-collision operator.  In $(\vp,\vpa)$ it
has the six-coefficient form \eqref{eq:FP}.

\paragraph{Energy confinement.}
With $\tau_E>0$ (\texttt{tau\_E}) and a source (\texttt{irf = -1} or \texttt{isource = -1}), every
species loses energy through
\begin{equation}
L_E(f)=\frac{1}{2\tau_E}\nabla_v\cdot(\bm v f)=\frac{1}{2\tau_E}\Big[3f+\vp\frac{\partial f}{\partial\vp}
+\vpa\frac{\partial f}{\partial\vpa}\Big],
\end{equation}
which conserves particles and gives $\mathrm dW_s/\mathrm dt=-W_s/\tau_E$ exactly (and
$\mathrm dP_{\parallel,s}/\mathrm dt=-P_{\parallel,s}/2\tau_E$).  It adds $3/2\tau_E$, $\vp/2\tau_E$ and
$\vpa/2\tau_E$ to $A$, $B$ and $C$.  Without a source $\tau_E$ is ignored (with a warning): the
plasma then relaxes at constant total energy.  At steady state the global balances are
\begin{equation}
P_\mathrm{RF}=\sum_s\frac{W_s}{\tau_E},\qquad
P_\mathrm{NBI}=\sum_s\frac{W_s}{\tau_E}+\frac{W_a}{\tau_s}.
\end{equation}
$L_E$ is a uniform contraction of velocity space: every particle loses energy at the same relative
rate, so tails whose collision time exceeds $\tau_E$ are depleted.  Momentum leaves the bulk only
through $L_E$, with the time $2\tau_E$; a beam injecting parallel momentum therefore spins the whole
plasma up, to drifts comparable with the ion thermal speed in the JET beam test.

\paragraph{Coulomb logarithms.}
NRL formulary: ion--ion (\texttt{coulomb\_log\_ab}, Section~\ref{sec:coulomb_log}), ion--electron and
electron--ion (\texttt{coulomb\_log\_ae}), and electron--electron
\begin{equation}
\ln\Lambda_{ee}=23.5-\ln\!\big(n_e^{1/2}T_e^{-5/4}\big)-\big[10^{-5}+(\ln T_e-2)^2/16\big]^{1/2}
\end{equation}
($n_e$ in cm$^{-3}$, $T_e$ in eV).  Each pair uses one value for both directions.  The temperatures
are updated every step: $T_\mathrm{eff}$ for the heated species, as in v3.2, and for a bulk species
the \emph{thermal} temperature $T_\mathrm{eff}-m u^2/3$, without the energy of its parallel drift
$u=P_\parallel/nm$.

\subsection{Species Grids}
\label{sec:ms_grids}

The heated species keeps the namelist grid.  Each bulk species $b$ gets the same grid type (power-law
$\vp$ grid with the same \texttt{ising}, \texttt{nsing}, \texttt{p\_grid}, rescaled; uniform $\vpa$) on
$\vp\le v_\mathrm{box}v_{t0}$, $|\vpa|\le v_\mathrm{box}v_{t0}$, where $v_{t0}=\sqrt{T_b/m_b}$ at the
initial temperature and $v_\mathrm{box}$ is \texttt{vbox\_b(ib)} (default 8).  The node numbers are
\texttt{nperp\_b(ib)}, \texttt{npar\_b(ib)} (default: those of the heated species).  The index
\texttt{ib} follows the bulk species: 1 for the electrons, 2, 3, \dots\ for the ions.  The initial
distribution is the Maxwellian at $T_b$, normalised to $n_b$ on the grid.

\paragraph{Choosing the boxes.}  A grid built for $T_0$ represents the species well for
$T/T_0$ between about $(4\Delta\vpa/v_{t0})^2$ and $(v_\mathrm{box}/5)^2$: beyond the upper bound the
walls truncate the distribution, below the lower one there are fewer than four cells per thermal
speed.  With 150 nodes this is $0.18$--$2.6$ for $v_\mathrm{box}=8$ and $0.4$--$5.8$ for
$v_\mathrm{box}=12$.  The coupled loop checks this at every step (Section~\ref{sec:ms_regrid}).
For runs whose temperatures change by larger factors, use the regridding
(\texttt{i\_regrid = -1}).  The ions occupy only one or two cells of the electron grid; the
collision coefficients are built so that this is not a problem (next subsection).

\subsection{Rosenbluth Potentials Across Grids}

The potentials of a background species are needed at the nodes of every target grid.  The v3.2
kernel works only on the background's own grid, so the coupled loop uses a Legendre expansion
(module \texttt{mod\_legendre\_pot}, \texttt{legendre\_pot.f90}):
\begin{enumerate}
\item $f_b(v,\xi)=\sum_{l=0}^Lf_l(v)P_l(\xi)$, $\xi=\vpa/v$, with $f_l$ sampled on a uniform radial grid
  (step $\Delta\vpa/4$ by default, \texttt{nr\_leg}) by Gauss--Legendre quadrature in $\xi$
  (\texttt{nxi\_leg} points, default $\max(64,4L+8)$) of a $4\times4$ Lagrange interpolant of $f_b$;
\item four cumulative radial moments of each $f_l$, integrated exactly for $f_l$ linear between radial
  nodes, from which the Legendre coefficients of $h$ and $g$ follow;
\item at any target node, $\psi$, $\phi$ and all their first and second derivatives are evaluated
  \emph{analytically} (no numerical differentiation), and $f_b$ itself replaces $\nabla^2\phi$.
\end{enumerate}
The degree is $L=$ \texttt{lmax\_leg} (default 16) for the heated species and $L_b=$
\texttt{lmax\_bulk} (default 4) for the bulk species.  $L_b\ge1$ is required for momentum
conservation; $L_b=4$ is accurate to $10^{-3}$ even for a bulk drift of 1.6 thermal speeds.  The
self-collisions of the heated species keep the v3.2 FFT kernel.

\subsection{Coupled Time Loop}
\label{sec:ms_loop}

Routine \texttt{timefp\_ms} (\texttt{TimeFP\_ms.f90}) replaces \texttt{timefp\_7pt\_nl} when
\texttt{i\_coupled = -1}.  At setup every species gets a fixed sparse pattern, quadrature weights and
its own PARDISO handle (the global matrix is block-diagonal and never assembled).  Each time step
then does:
\begin{enumerate}
\item \emph{Regridding} of any bulk species whose grid no longer fits (\texttt{i\_regrid = -1});
  $f^*=f^n$.
\item \emph{Picard pass} $p=0,\dots,n_\mathrm P$ (\texttt{n\_picard}):
  \begin{enumerate}
  \item moments, temperatures, drifts and Coulomb logarithms of every $f^*$; grid check;
  \item Legendre representation of every $f^*$ on its own grid;
  \item the coefficients of every pair $(s\leftarrow j)$ on the grid of $s$, plus $Q_\mathrm{RF}$ (heated
    species) and $L_E$ (every species);
  \item one $\theta$-step per species: $(I-\theta\Delta t\,L_s)f_s^{n+1}=(I+(1-\theta)\Delta t\,L_s)f_s^n
    +\Delta t\,S_s$;
  \item the moment projection (below); then $f^*=\theta f^{n+1}+(1-\theta)f^n$ for the next pass.
  \end{enumerate}
\item RF, confinement, source and particle-loss powers on the corrected state; diagnostics.
\end{enumerate}
With \texttt{n\_picard = 0} the potentials are lagged by one step (the iteration scheme of v3.2
extended to all pairs); \texttt{n\_picard = 1} evaluates them at the time level of the
$\theta$-scheme at twice the cost and is recommended.  The two converge to the same steady state as
$\Delta t\to0$ (0.2--0.3\% apart in $T_e$ at $\Delta t=25$\,ms, $10^{-5}$ at 10\,ms in the RF test).

\paragraph{Moment projection (\texttt{i\_ms\_cons}).}  The discrete pair operators do not conserve
particles, energy and momentum exactly: their $\sim10^{-4}$ errors are multiplied by the large
electron self-collision frequency, the ions are unresolved on the electron grid, and lagged
potentials unbalance the pairs.  Uncorrected, an ion--electron equilibration test loses 72\% of its
electrons in 3\,s.  After each step the code measures the violations of every species from the rates
of every pair (an inter-species defect shared $n$-weighted for energy and $\rho$-weighted for
momentum) and removes them with the implicit responses of the species to three operators that
change only $n$, $W$ or $P_\parallel$ of a Maxwellian (growth, dilation, shift).  The amplitudes are
chosen so that the conservation laws hold for the \emph{corrected} state, one linear system of size
$3n_s$.  \texttt{i\_ms\_cons = 2} (default) conserves all particle numbers, the total energy and the
total parallel momentum to $10^{-10}$ or better; it modifies the operator by about $10^{-4}$ of its
own rate.  Values 0 (none), 1 (particles only) and $-2$ (explicit amplitudes) exist for comparison
and should not be used for physics runs.

\paragraph{Cost.}  On three $150\times150$ grids (8 threads) a coupled step takes about 0.14\,s with
lagged potentials and 0.28\,s with one Picard iteration; the three linear solves dominate.  The
end of the run prints a wall-clock profile of the step.

\subsection{Sources and Losses in the Coupled Loop}

RF heating (\texttt{irf = -1}) and the beam (\texttt{isource = -1}) act on the heated species as in
v3.2; $L_E$ acts on every species.  For a beam run, \texttt{istart = 0} starts from $f_a=0$; the
self-collisions of the beam are skipped while its density is below 5\% of the final one, as in v3.2.
The heated species may also start from a Maxwellian at \texttt{ta\_init} (eV), drifting at
\texttt{ua\_init} (m/s).  The source, loss and confinement rates are evaluated on the corrected
state, so the energy, beam-particle and momentum balances close at every step to round-off:
\begin{equation}
\frac{W^{n+1}_\mathrm{tot}-W^n_\mathrm{tot}}{\Delta t}=P_\mathrm{RF}+P_\mathrm{NBI}+P_\mathrm{loss}+\sum_sP_{E,s}.
\end{equation}

\paragraph{Second-harmonic RF.}  At $\omega=2\omega_c$ the absorbed power grows with the energy of the
heated species, like the loss $\sum_sW_s/\tau_E$.  With evolving backgrounds and a common $\tau_E$
there is then no equilibrium at moderate temperature: the plasma cools or runs away to MeV energies.
Fundamental heating, or a beam, has a stable equilibrium.

\subsection{Grid Checks and Regridding}
\label{sec:ms_regrid}

\paragraph{Grid check.}  At every step the coupled loop warns once per bulk species when its wall is
closer than $5v_t$ to its drift (thermal $v_t$), or when $v_t<4\Delta\vpa$, and stops with a message if
a temperature becomes non-positive.

\paragraph{Regridding (\texttt{i\_regrid = -1}).}  At the start of a step a bulk species is
regridded when its wall distance from the drift falls below $v_\mathrm{box}v_t/\phi$ (heating or
drift) or its half-width exceeds $\phi\,v_\mathrm{box}v_t$ (cooling), with
$\phi=$ \texttt{regrid\_fac} (default 1.3): one regrid each time $T$ changes by about $\phi^2=1.7$.
The new grid keeps the node numbers, has the half-width $v_\mathrm{box}v_t$ and its $\vpa$ range
centred on the drift.  The distribution is transferred as the drifting Maxwellian with the same
$n$, $u$, $T$ plus the interpolated deviation of the old one, and $n$, $W$ and $P_\parallel$ are then
restored exactly, so the conservation laws and step balances are continuous across a regrid.  The
species' solver data are rebuilt.  In a tenfold ion heating the regridded run agrees with a fixed
fine grid to 6\,eV, where fixed grids sized for the initial temperature are 340\,eV off.  The heated
species is not regridded: size its namelist grid for the whole run.

\subsection{Verification Modes and Checks}
\label{sec:ms_checks}

\begin{description}
\item[\texttt{i\_bulk\_frozen = -1}] (coupled loop) holds the bulk species at their initial
  Maxwellians, without $L_E$ and without projection: the coupled loop then solves the v3.2 problem
  with its own operators, and must reproduce a v3.2 run (\texttt{i\_coupled = 0}) to $10^{-3}$ of $f$.
\item[\texttt{i\_ms\_check = -1}] (with \texttt{i\_multi = -1}) compares, at start, the coefficients of
  every pair built from the Legendre potentials of the initial Maxwellians with the analytic
  Maxwellian operator \texttt{cblin}, and, at the end of a v3.2 nonlinear run, the Legendre potentials
  of the final $f$ with the FFT kernel.  Results in \texttt{ms\_check.txt} and maps in
  \texttt{ms\_cblin\_*.dat}.
\item[\texttt{i\_sc00 = -1}] uses the exact $c_{00}=-\kappa f$ in the v3.2 nonlinear self-collisions
  instead of the finite-difference Laplacian of $\phi$.  It is a study option; the default is
  recommended (the finite-difference value is consistent with the other discrete coefficients).
\end{description}

\subsection{Practical Guidelines}

\begin{itemize}
\item Start from a working v3.2 namelist; add \texttt{i\_multi = -1}, \texttt{i\_coupled = -1}
  (with \texttt{isc = -1}, \texttt{iold = 0}), \texttt{n\_picard = 1}, and \texttt{tau\_E} when there is a
  source.  Keep \texttt{i\_ms\_cons = 2}.
\item Bulk boxes: \texttt{vbox\_b} 8 for the electrons and 12 for the ions suit temperatures within a
  factor of about two of the initial ones; otherwise set \texttt{i\_regrid = -1}.
\item $\tau_E$: the steady state has $\sum_sW_s\simeq P\,\tau_E$, so choose $\tau_E$ from the expected power
  and stored energy.  Check the end-of-run temperatures against the grids (the grid warnings).
\item Time step: 5--25\,ms have been used.  Lagged and Picard steady states differ by 0.2--0.3\% at
  25\,ms and agree to $10^{-5}$ at 10\,ms.  Use smaller steps at the start of a beam fill-up or of a
  fast transient.
\item There is no automatic steady-state stop in the coupled loop (\texttt{i\_ss\_check} is ignored):
  set \texttt{ntimes} for the expected relaxation (several $\tau_E$, or $2\tau_E$ for the rotation of a
  beam case).
\item Runs with several threads are not bit-reproducible; use \texttt{n\_threads = 1} for byte
  comparisons.
\end{itemize}

\subsection{Limitations and Untested Configurations}

\begin{itemize}
\item \emph{Verified range}: one heated species, electrons and \emph{one} bulk ion species, all with
  $Z=1$; RF or NBI.  More bulk ion species (up to ten) and $Z>1$ are allowed but untested; run a
  multi-ion check before using them for D--T--$\alpha$ or impurity cases.
\item RF and NBI together, restarts (\texttt{iold = -1}) and the steady-state solver are not available
  or not tested in the coupled loop.
\item Non-relativistic electrons; momentum lost only through $L_E$ (beam cases over-rotate); no
  fixed-power RF option; the heated species is not regridded.
\end{itemize}


% ============================================================
\section{Namelist Reference}
\label{sec:namelist}
% ============================================================

The code reads its input from standard input (\texttt{stdin}) as a Fortran
\texttt{NAMELIST} group named \texttt{INPUT}.  A typical invocation is:
\begin{lstlisting}
FP2D_QLRF_NL.exe < jet_beam_case1.dat
\end{lstlisting}
The namelist file may omit any variable that has a default value; the code
will use the default.

\bigskip
The parameters are grouped thematically below.

% -------------------------------------------------------------------
\subsection{General and Output Control}

\begin{longtable}{>{\ttfamily}p{3.2cm} p{1.5cm} p{9cm}}
\toprule
\normalfont Parameter & Default & Description \\
\midrule
\endhead
\bottomrule
\endlastfoot

casename & \texttt{''} &
  Character string (max.\ 64 chars) appended to every output filename and
  plot title, e.g.\ \texttt{'JET-beam7-TD0-NLSC'}.  Leave empty for
  unnamed runs. \\[4pt]

iplot\_pow & $-1$ &
  Controls writing of power-vs-time output files.
  $-1$: write all power files.
  $0$: skip (useful when only energy and density diagnostics are needed). \\[4pt]

iplot\_mom & $0$ &
  Controls writing of momentum-transfer-rate files.
  $-1$: write all momentum files.
  $0$: skip (default; momentum files are large). \\[4pt]

n\_snap & $0$ &
  VDF snapshots during a time-dependent run.
  $0$: none.
  $N > 0$: write $f$ every $N$ steps to \texttt{vdf\_snap\_<step>.txt}
  (Section~\ref{sec:snap}).  Both time solvers. \\[4pt]

notxt & $0$ &
  Controls how many \texttt{.txt} output files are written (the \texttt{.dat}
  restart/kernel files are never affected).
  $0$: write all \texttt{.txt} outputs (normal).
  $1$: suppress all \texttt{.txt} outputs.
  $2$: write only the minimal test/diagnostic set, namely
  \texttt{Ekin\_perp}, \texttt{Ekin\_perp\_at\_vpar0}, \texttt{fout\_at\_vpar0},
  \texttt{Teff\_vs\_time}, and \texttt{anisotropy\_vs\_time}.
  Suppressed files are routed to the OS null device, so the run is otherwise
  unchanged (all stdout diagnostics still print). \\[4pt]

\end{longtable}

% -------------------------------------------------------------------
\subsection{Velocity-Space Grid}

\begin{longtable}{>{\ttfamily}p{3.2cm} p{1.5cm} p{9cm}}
\toprule
\normalfont Parameter & Default & Description \\
\midrule
\endhead
\bottomrule
\endlastfoot

nperp & -- &
  Number of grid points in $\vp$ (perpendicular direction). \\[4pt]

npar & -- &
  Number of grid points in $\vpa$ (parallel direction). \\[4pt]

vperp\_min & -- &
  Minimum perpendicular velocity (m/s).  Should be 0 or very small. \\[4pt]

vperp\_max & -- &
  Maximum perpendicular velocity (m/s).  Should be at least $3$--$5\times$
  the beam or thermal velocity. \\[4pt]

vpar\_min & -- &
  Minimum parallel velocity (m/s).  Must be negative. \\[4pt]

vpar\_max & -- &
  Maximum parallel velocity (m/s).  Must be positive and equal to
  $|$\texttt{vpar\_min}$|$ for symmetric grids. \\[4pt]

ising & $0$ &
  Grid spacing mode in $\vp$.
  $0$: uniform.
  $-1$: two-domain (fine for $\vp \le \texttt{vbound}$, coarser above).
  $+1$: power-law spacing (quadratic by default, see \texttt{p\_grid}), higher
  resolution near $\vp = 0$. \\[4pt]

p\_grid & $2$ &
  Exponent of the power-law grid when \texttt{ising = +1}:
  ${\vp}_i = \texttt{vperp\_min} + \texttt{vperp\_max}\,
  \bigl[(i-\tfrac12)/(\texttt{nperp}-1)\bigr]^{p}$.
  $p = 2$ is the original quadratic grid (reproduced to the bit); a larger
  $p$ moves points from the tail into the thermal bulk and coarsens the
  spacing at the $\vp$ wall, which raises the wall P\'eclet number.
  Ignored for \texttt{ising} $\ne +1$. \\[4pt]

nsing & -- &
  Number of fine grid points in the inner domain when \texttt{ising = -1}.
  Must satisfy $1 \le \texttt{nsing} < \texttt{nperp}$. \\[4pt]

vbound & -- &
  Boundary between fine and coarse regions in $\vp$ when \texttt{ising = -1}
  (m/s). \\[4pt]

\end{longtable}

% -------------------------------------------------------------------
\subsection{Plasma Parameters}

\begin{longtable}{>{\ttfamily}p{3.2cm} p{1.5cm} p{9cm}}
\toprule
\normalfont Parameter & Default & Description \\
\midrule
\endhead
\bottomrule
\endlastfoot

nbulk & -- &
  Number of background species (including electrons).
  Species 1 is always electrons.  Up to 10 species are supported. \\[4pt]

t(1:nbulk) & -- &
  Temperatures of background species (eV).
  \texttt{t(1)}: electron temperature $T_e$.
  \texttt{t(2)}: temperature of first ion species, etc. \\[4pt]

aa & -- &
  Mass number of the minority (test) ion species. \\[4pt]

za & -- &
  Charge number (atomic number) of the minority ion species. \\[4pt]

ab(1:9) & -- &
  Mass numbers of background ion species 2, 3, \ldots\ (species 1 is
  always electrons with $m_e/m_p$). \\[4pt]

zb(1:9) & -- &
  Charge numbers of background ion species 2, 3, \ldots \\[4pt]

ne & -- &
  Electron density (m$^{-3}$). \\[4pt]

xpart & -- &
  Concentration of minority ion species relative to $n_e$,
  i.e.\ $n_a = \texttt{xpart} \times n_e$. \\[4pt]

xb(1:9) & -- &
  Concentrations of background ion species 2, 3, \ldots\ relative to $n_e$.
  Quasi-neutrality is assumed: $\sum_b Z_b\,n_b = n_e$. \\[4pt]

\end{longtable}

% -------------------------------------------------------------------
\subsection{Beam Source}

\begin{longtable}{>{\ttfamily}p{3.2cm} p{1.5cm} p{9cm}}
\toprule
\normalfont Parameter & Default & Description \\
\midrule
\endhead
\bottomrule
\endlastfoot

isource & -- &
  Beam source flag.
  $-1$: include NBI beam source and particle losses.
  $0$: no source; distribution is sourceless (total particles conserved). \\[4pt]

beam\_ekin & -- &
  Beam injection energy (keV). \\[4pt]

beam\_angle\_deg & -- &
  Injection angle (degrees) measured from the magnetic field direction.
  $0^\circ$: purely parallel injection.
  $90^\circ$: purely perpendicular injection. \\[4pt]

beam\_dvperp & -- &
  Width $\sigma_\perp$ of the Gaussian approximation to the delta function
  in $\vp$ (m/s).  A value of $\sim 5$--$10\%$ of $\vp^\mathrm{beam}$ is
  typical; must be resolved by the grid. \\[4pt]

beam\_dvpar & -- &
  Width $\sigma_\parallel$ of the Gaussian in $\vpa$ (m/s). \\[4pt]

taus & -- &
  Slowing-down time $\tau_s$ (s).  Sets the source amplitude
  $S_0 = n_a / \tau_s$ and the loss rate. \\[4pt]

\end{longtable}

% -------------------------------------------------------------------
\subsection{RF Heating}

\begin{longtable}{>{\ttfamily}p{3.2cm} p{1.5cm} p{9cm}}
\toprule
\normalfont Parameter & Default & Description \\
\midrule
\endhead
\bottomrule
\endlastfoot

irf & -- &
  RF flag.
  $-1$: include quasi-linear RF operator.
  $0$ or else: no RF. \\[4pt]

frek & -- &
  Wave frequency $f$ (Hz). \\[4pt]

b0 & -- &
  Static magnetic field $B_0$ (T). \\[4pt]

nharm & -- &
  Resonance harmonic number $n$ (e.g.\ 1 for fundamental, 2 for second
  harmonic). \\[4pt]

kpar & -- &
  Parallel wave vector $k_\parallel$ (rad/m). \\[4pt]

kperp & -- &
  Perpendicular wave vector $k_\perp$ (complex, rad/m). \\[4pt]

eplus & -- &
  Left-hand polarisation amplitude $E_+$ (complex, V/m). \\[4pt]

emin & -- &
  Right-hand polarisation amplitude $E_-$ (complex, V/m). \\[4pt]

delta\_RF & -- &
  Width of the Gaussian broadening of the resonance in $\vpa$ (m/s).
  Must be wide enough to resolve the resonance on the grid. \\[4pt]

\end{longtable}

% -------------------------------------------------------------------
\subsection{Time Integration and Solver Control}

\begin{longtable}{>{\ttfamily}p{3.2cm} p{1.5cm} p{9cm}}
\toprule
\normalfont Parameter & Default & Description \\
\midrule
\endhead
\bottomrule
\endlastfoot

ntimes(1:3) & \texttt{0,0,0} &
  Number of time steps per phase.
  \texttt{ntimes(1) = 0}: steady-state run.
  \texttt{ntimes(1) > 0}: time-dependent run.
  Up to three phases with different step sizes are supported; phases with
  \texttt{ntimes(k) = 0} are skipped.
  Old-format scalar values (\texttt{ntimes = N}) are backward-compatible and
  fill \texttt{ntimes(1)}. \\[4pt]

timestep(1:3) & \texttt{0,0,0} &
  Time step size per phase (s).
  \texttt{timestep(k)} is used for phase~$k$.
  Backward-compatible with scalar input. \\[4pt]

icn & -- &
  Time-differencing scheme.
  $-1$: Crank--Nicolson ($\theta = 0.5$).
  $+1$: intermediate scheme ($\theta = 0.75$).
  else: fully implicit ($\theta = 1.0$). \\[4pt]

iold & -- &
  Initial condition source.
  $0$: fresh start (use \texttt{istart} to select the initial condition).
  $-1$: restart from a previous run by reading \texttt{xout.dat}. \\[4pt]

istart & $1$ &
  Initial distribution for fresh time-dependent runs (\texttt{iold = 0}).
  $0$: $f^0 = 0$ (beam-driven only; requires \texttt{isource = -1}).
  $1$: Stix Maxwellian.
  $2$: Steady-state solution without self-collisions.
  $3$: Steady-state solution with Maxwellian self-collisions. \\[4pt]

isc & -- &
  Self-collision treatment.
  $0$: no self-collisions (linear FP equation).
  $1$: Maxwellian background at fixed Stix temperature.
  $2$: Maxwellian background at evolving effective temperature $T_\mathrm{eff}(t)$
       (requires TD run).
  $3$: Maxwellian background at the density-characteristic temperature $T_n(t)$
       (requires TD run; see \texttt{core\_frac}).
  $-1$: fully nonlinear self-collision operator (requires TD run). \\[4pt]

core\_frac & $3.8\!\times\!10^{-3}$ &
  Core fraction defining the $T_n$ log-slope fit for \texttt{isc = 3}: only grid
  cells with $f > \texttt{core\_frac}\cdot\max f$ enter the $\ln f$-vs-$v^2$ fit.
  Smaller values weight the cold bulk more heavily.  Ignored unless
  \texttt{isc = 3}; the appropriate value is scenario-dependent. \\[4pt]

i\_upwind & $0$ &
  Advection stabilization of the perpendicular drag term $B\,\partial f/\partial\vp$
  (Section~\ref{sec:upwind}).
  $0$: central Fornberg weights (default).
  $-1$: P\'eclet-hybrid upwinding --- in cells where the cell-P\'eclet number
  $|B|\,\Delta\vp/D > 2$ the drag term uses a one-sided (first-order) upwind
  difference, restoring dissipativity of the advection-dominated outer boundary
  layer and removing the Crank--Nicolson high-$\vp$ instability.
  $1$: P\'eclet-hybrid \emph{power-law} scheme (Patankar 1980) --- in the same
  advection-dominated cells the combined drag$+$diffusion operator is
  discretised with a 3-point control-volume power-law stencil, which removes the
  instability with roughly half the tail numerical diffusion of the first-order
  upwind ($-1$).  Both leave the high-order 7-point Fornberg stencil in the
  smooth bulk. \\[4pt]

i\_lu\_reuse & $-1$ &
  Linear solves of the time steps.
  $-1$: reuse the LU factors of the last factorisation as the preconditioner of a
  CGS iteration, and refactorise only when its result is rejected
  (Section~\ref{sec:lureuse}); 2.7--3.4$\times$ faster runs.
  $0$: refactorise the matrix at every step (the former behaviour). \\[4pt]

lu\_reuse\_tol & $10^{-10}$ &
  Residual bound for \texttt{i\_lu\_reuse = -1}: a reused step is accepted if
  $\|b - A x\|_\infty/\|b\|_\infty \le \max(\texttt{lu\_reuse\_tol},\,
  10\,r_\mathrm{direct})$, with $r_\mathrm{direct}$ the residual of the latest
  direct solve.  Values below ${\sim}10^{-11}$ make every step refactorise. \\[4pt]

n\_threads & $0$ &
  Number of threads for MKL (PARDISO) and the OpenMP loops (Release build).
  $0$: \texttt{MKL\_NUM\_THREADS} / \texttt{OMP\_NUM\_THREADS} if set, otherwise
  16.  $n>0$: $n$ threads for both, overriding the environment.  $1$ is the
  sequential code: results identical to the sequential (Debug) build and
  reproducible from run to run; with more threads they differ in the last
  digits.  No effect in the Debug build (Section~\ref{sec:threads}). \\[4pt]

\end{longtable}

% -------------------------------------------------------------------
\subsection{Steady-State Convergence Detection}

\begin{longtable}{>{\ttfamily}p{3.2cm} p{1.5cm} p{9cm}}
\toprule
\normalfont Parameter & Default & Description \\
\midrule
\endhead
\bottomrule
\endlastfoot

i\_ss\_check & $0$ &
  Enable automatic convergence detection (Section~\ref{sec:ssconv}).
  $0$: disabled (run for the full \texttt{ntimes} steps).
  $-1$: stop early when all five criteria are met; also writes
  \texttt{conv\_eps\_vs\_time.txt} and \texttt{conv\_diag\_vs\_time.csv}. \\[4pt]

n\_ss\_window & $50$ &
  Window width $W$ (number of time steps).  The convergence test is evaluated
  once every $W$ steps, comparing the current state with the one stored at the
  previous evaluation and dividing by the physical duration the window spans.  A
  wider window is less sensitive to short-term fluctuation and proportionally
  cheaper, but coarsens the time resolution of the diagnostic and delays the
  first verdict (which comes at step $2W$).  $W = 1$ evaluates every step. \\[4pt]

ss\_tol\_eps & $10^{-3}$ &
  Tolerance on $\varepsilon/\nu_\mathrm{ref}$, the Jacobian-weighted bulk
  convergence rate, Eq.~\eqref{eq:epsconv}.  Dimensionless (a fraction of the
  reference rate) and independent of \texttt{timestep}. \\[4pt]

ss\_tol\_tail & $10^{-2}$ &
  Tolerance on $\varepsilon_\mathrm{tail}/\nu_\mathrm{ref}$, the tail-weighted
  rate.  Looser than \texttt{ss\_tol\_eps} by default because the tail is the
  noisier, least-populated part of the distribution; tighten it when the physics
  of interest lives in the tail. \\[4pt]

ss\_tol\_moment & $10^{-3}$ &
  Tolerance on the relative moment drift rates over the window.  Which moments
  are tested depends on \texttt{i\_conv\_shape}: with $0$ the density, parallel
  flow and energy; with $-1$ the intensive moments only, i.e.\ the parallel flow
  and the mean energy per particle $E/n$. \\[4pt]

i\_conv\_shape & $0$ &
  Selects \emph{which family of rates the tolerances are applied to}
  (Section~\ref{sec:ssconv}).  Both families are always computed and written out;
  this flag changes only which ones are tested.

  $0$ --- \emph{amplitude}.  The criterion uses $\varepsilon$ and
  $\varepsilon_\mathrm{tail}$, Eq.~\eqref{eq:epsconv}, together with the drift
  rates of $n$, $u_\parallel$ and $E$.  Five conditions in all.  This is the
  correct choice when the run reaches a genuine steady state, e.g.\ when a beam
  source balances the boundary losses.

  $-1$ --- \emph{shape}.  $f$ is normalised to unit norm before differencing,
  Eq.~\eqref{eq:epsshape}, so the criterion uses $\varepsilon_\mathrm{shape}$ and
  $\varepsilon_\mathrm{tail,shape}$ together with the drift rates of the
  \emph{intensive} moments $u_\parallel$ and $E/n$; the density drift is dropped,
  being pure amplitude.  Four conditions in all.  Required for a
  \emph{sourceless} run: particles absorbed at the Dirichlet boundaries are never
  replaced, so $f$ decays as a frozen shape with falling amplitude and the
  amplitude rates floor at the drain rate, which no tolerance below that rate can
  ever clear. \\[4pt]

\end{longtable}

Convergence requires \emph{all} the selected conditions to be satisfied
simultaneously at one evaluation.  If a run never stops, compare
$\varepsilon$ with $\varepsilon_\mathrm{shape}$ in
\texttt{conv\_eps\_vs\_time.txt}: if the former has flattened while the latter is
still falling, the run is draining and needs \texttt{i\_conv\_shape = -1}.


% -------------------------------------------------------------------
\subsection{Multispecies (Version 4.0)}
\label{sec:nml_ms}

All defaults reproduce version 3.2.  Arrays indexed by \texttt{ib} follow the bulk species
(1 = electrons, 2, 3, \dots\ = bulk ions).

\begin{longtable}{>{\ttfamily}p{3.2cm} p{1.5cm} p{9cm}}
\toprule
\normalfont Parameter & Default & Description \\
\midrule
\endhead
\bottomrule
\endlastfoot

i\_multi & $0$ &
  $0$: single species (v3.2).  $-1$: one container, grid and distribution per species
  (Section~\ref{sec:ms}). \\[4pt]

tau\_E & $0$ &
  Energy-confinement time (s) of $L_E$, applied to every species (to the heated one only when
  \texttt{i\_coupled = 0}), and only with a source (\texttt{irf = -1} or \texttt{isource = -1});
  ignored with a warning otherwise.  $0$: no loss. \\[4pt]

nperp\_b(ib), npar\_b(ib) & $0$ &
  Node numbers of the grid of bulk species \texttt{ib}.  $0$: those of the heated species. \\[4pt]

vbox\_b(ib) & $0$ &
  Box of bulk species \texttt{ib}, in thermal speeds at its initial temperature.  $0$: 8. \\[4pt]

i\_coupled & $0$ &
  $0$: backgrounds frozen (v3.2 physics).  $-1$: coupled evolution of all species
  (\texttt{timefp\_ms}); needs \texttt{i\_multi = -1}, \texttt{isc = -1}, \texttt{iold = 0} and
  \texttt{ntimes(1) > 0}. \\[4pt]

n\_picard & $0$ &
  Picard iterations per coupled step.  $0$: potentials lagged by one step; $1$ (recommended):
  potentials at the $\theta$ level, twice the cost. \\[4pt]

i\_ms\_cons & $2$ &
  Conservation of the coupled step (implicit moment projection).  $2$: particles of each species,
  total energy and total momentum.  $1$: particles only; $0$: none; $-2$: explicit amplitudes.
  Values other than 2 are for comparison only. \\[4pt]

lmax\_leg & $16$ &
  Legendre degree of the heated species' potentials ($0$--$48$). \\[4pt]

lmax\_bulk & $4$ &
  Legendre degree of the bulk species' potentials ($0$--$48$; $\ge1$ for momentum
  conservation). \\[4pt]

nr\_leg, nxi\_leg & $0$, $0$ &
  Radial intervals and Gauss points in $\xi$ of the Legendre expansion.  $0$: automatic (radial
  step $\Delta\vpa/4$; $\max(64,4L+8)$ points). \\[4pt]

ta\_init, ua\_init & $0$, $0$ &
  Coupled run: initial Maxwellian of the heated species at \texttt{ta\_init} (eV), drifting at
  \texttt{ua\_init} (m/s).  \texttt{ta\_init = 0}: the start selected by \texttt{istart}. \\[4pt]

i\_bulk\_frozen & $0$ &
  $-1$: bulk species held at their initial Maxwellians in the coupled loop (verification of the
  v3.2 limit). \\[4pt]

i\_regrid & $0$ &
  $-1$: regrid a bulk species when its grid no longer fits its thermal speed and drift
  (Section~\ref{sec:ms_regrid}). \\[4pt]

regrid\_fac & $1.3$ &
  Regridding tolerance $\phi>1$: regrid when the wall distance falls below
  $v_\mathrm{box}/\phi$ or the half-width exceeds $\phi\,v_\mathrm{box}$ thermal speeds. \\[4pt]

i\_ms\_check & $0$ &
  $-1$: checks of the Legendre potentials (with \texttt{i\_multi = -1};
  Section~\ref{sec:ms_checks}). \\[4pt]

i\_sc00 & $0$ &
  $-1$: exact $c_{00}$ in the v3.2 nonlinear self-collisions (study option). \\[4pt]

\end{longtable}


% ============================================================
\section{Output Files}
% ============================================================

All output files are written to the working directory.  When \texttt{casename}
is non-empty, the string \texttt{-casename} is inserted before the file
extension.

\subsection{Final Distribution}

\begin{description}
  \item[\texttt{fout.txt}] Final VDF as a list of $(\vp, \vpa, f)$ triples.
  \item[\texttt{xout.dat}] Restart file: one header line, \texttt{time nstep}
        (final time and total number of steps, including those of any run it
        was itself restarted from), followed by the packed 1D solution vector.
        Used with \texttt{iold = -1}.  Older files whose header holds only the
        time are still accepted; the step count then starts from 0.
  \item[\texttt{fstix.dat}] Stix Maxwellian reference distribution.
  \item[\texttt{Ekin.txt}, \texttt{Ekin\_perp.txt}, \texttt{Ekin\_par.txt},
        \texttt{Ekin\_perp\_at\_vpar0.txt}] Kinetic-energy density of the final
        $f$ (total, perpendicular and parallel parts), and its cut at
        $\vpa = 0$.
  \item[\texttt{fsc\_maxw.txt}, \texttt{fsc\_maxw\_at\_vpar0.txt}] Maxwellian
        background of the self-collision operator, at $T_\mathrm{eff}$
        (\texttt{isc = 2}) or $T_n$ (\texttt{isc = 3}).
  \item[\texttt{beam.txt}] NBI source $S(\vp,\vpa)$ (\texttt{isource = -1}).
\end{description}

\subsection{VDF Snapshots}
\label{sec:snap}

With \texttt{n\_snap = N > 0}, $f$ is written every $N$ steps to
\texttt{vdf\_snap\_<step>.txt}, in the same layout as \texttt{fout.txt} with a
\texttt{\# time =} header line.  \texttt{<step>} is the \emph{total} step count
read from the \texttt{xout.dat} header, so a restart continues the previous
run's numbering and cadence instead of overwriting its snapshots.  Snapshots
hold the raw solution (the end-of-run renormalisation of sourceless runs is not
applied) and are not restart files.  \texttt{fp2d\_plot.py --movie} turns them
into an animation (Section~\ref{sec:plotting}).

\subsection{1D Cuts of the VDF}

\texttt{fout\_at\_vpar0.txt}, \texttt{fout\_at\_vperp0.txt}: values of $f$
along the lines $\vpa = 0$ and $\vp = 0$.  (The cuts at $\vpmax$ and $\vpamax$,
and those of the Stix Maxwellian, are no longer written.)

\subsection{Time Traces}

Written at every time step during the simulation (except the convergence files,
which are written once per \texttt{n\_ss\_window}-step window):

\begin{description}
  \item[\texttt{density\_vs\_time.txt}] Columns: time (s), $n(t)$ (m$^{-3}$).
  \item[\texttt{energy\_vs\_time.txt}] Columns: time, $E_\mathrm{tot}$ (keV), $E_\perp$ (keV).
  \item[\texttt{anisotropy\_vs\_time.txt}] Columns: time, perpendicular anisotropy (\%).
  \item[\texttt{Teff\_vs\_time.txt}] Columns: time, $T_\mathrm{eff}$ (keV).
  \item[\texttt{Tn\_vs\_time.txt}] Columns: time, $T_n$ (keV) -- density-characteristic
        (cold-bulk) temperature; written for all runs.
  \item[\texttt{conv\_eps\_vs\_time.txt}] Columns: time, $\varepsilon$,
        $\varepsilon_\mathrm{tail}$, $\varepsilon_\mathrm{shape}$,
        $\varepsilon_\mathrm{tail,shape}$ (s$^{-1}$) -- the convergence rates,
        Eqs.~\eqref{eq:epsconv} and~\eqref{eq:epsshape}.  Written when
        \texttt{i\_ss\_check = -1}, one row per \texttt{n\_ss\_window}-step window.
  \item[\texttt{conv\_diag\_vs\_time.csv}] Full convergence history: the four rates
        above, their ratios to $\nu_\mathrm{ref}$, the moments $n$, $u_\parallel$,
        $E$, $E/n$, all four drift rates and the unweighted axis-row rate.
  \item[\texttt{power\_coll\_tot\_vs\_time.txt}] Total collisional power density (MW/m$^3$).
  \item[\texttt{power\_coll\_e\_vs\_time.txt}] Electron collisional power density.
  \item[\texttt{power\_coll\_ion$N$\_vs\_time.txt}] Ion $N$ collisional power density.
  \item[\texttt{power\_coll\_self\_vs\_time.txt}] Self-collision power density
        (4 columns: time, total, $P_\mathrm{SC}^\perp$, $P_\mathrm{SC}^\parallel$).
  \item[\texttt{power\_NBI\_vs\_time.txt}] Beam injection and particle loss power.
  \item[\texttt{power\_RF\_vs\_time.txt}] RF power density.
  \item[\texttt{coulomb\_log\_vs\_time.txt}] Coulomb logarithm for background ion species.
  \item[\texttt{coulomb\_log\_self\_vs\_time.txt}] Self-collision Coulomb logarithm.
  \item[\texttt{density\_terms\_vs\_time.txt}] Particle rate $\mathrm{d}n/\mathrm{d}t$
        carried by each term of the operator (m$^{-3}$\,s$^{-1}$): time, total,
        collisions with each background species (electrons first),
        self-collisions, RF, source, losses.  Each term is the velocity
        integral of its own part of $\mathcal{L}f$; every term is a divergence,
        so a non-zero value means particles leave through the Dirichlet walls
        or the term is not conservative in its discrete form.  Switched off
        since version~3.2 (the lines are kept, commented out, in both time
        solvers).
  \item[\texttt{momentum\_coll\_*\_vs\_time.txt}] Momentum transfer rates
        (2 columns per file: $\perp$ and $\parallel$ components), written when
        \texttt{iplot\_mom = -1}.  The RF, NBI and self-collision rates go
        to \texttt{momentum\_RF}, \texttt{momentum\_NBI} and
        \texttt{momentum\_coll\_self}.
  \item[\texttt{tau\_coll\_vs\_time.txt}] Columns: time, effective
        $\tau_{ii}$, $\tau_{ie}$ (s).
  \item[\texttt{tau\_rf\_vs\_time.txt}] Columns: time, RF tail-formation
        time (s); written when \texttt{irf = -1}.
\end{description}

\subsection{Other Files}

\begin{description}
  \item[\texttt{phi\_kern-<casename>.dat}] Cache of the compressed
        $\phi$-kernel of the nonlinear self-collision operator
        (\texttt{isc = -1}); reused by a later run on the same grid, rebuilt
        otherwise.  Not a diagnostic; may be deleted.
  \item[\texttt{RF\_dirac.txt}] Parallel resonance function of the RF
        operator (\texttt{irf = -1}, \texttt{notxt = 0}).  Not routed through
        \texttt{casename}: every run overwrites it.
\end{description}

\subsection{End-of-Run Diagnostics (stdout)}
\label{sec:endofrun}

\paragraph{Balances.}  The power and particle-density balances are printed to
standard output at the end of every run, the momentum ($\perp$ and
$\parallel$) balance only when \texttt{iplot\_mom = -1}.  Values are given in
MW/m$^3$ for power, N/m$^3$ for momentum, and s$^{-1}$ for particle rates.
(The former switch \texttt{idiag}, which selected them, has been removed; a
namelist that still sets it is rejected at read time.)

\paragraph{Domain adequacy check.}  The Dirichlet walls impose $f=0$, whereas
the true solution there is $f(\text{wall})/f(\text{peak}) =
\exp(-V^2/2v_\mathrm{th}^2)$, with $V$ the wall velocity.  That ratio is a floor
on the error which grid refinement does not reduce.  At the end of every run the
code takes the \emph{final} $T_\mathrm{eff}$ (an RF or beam case may suit its
initial Maxwellian and still be too small for the heated distribution),
computes $v_\mathrm{th}=\sqrt{T/m}$, and prints, for $\vpmax$ and for
$\min(|\vpamin|,|\vpamax|)$, the ratio $f(\text{wall})/f(\text{peak})$ and a
verdict \texttt{OK} or \texttt{TOO SMALL} against a target of $10^{-12}$
($|v|_{\max}\approx 7.4\,v_\mathrm{th}$; this gave a shape error of
${\sim}2\times10^{-8}$, against ${\sim}3\times10^{-5}$ at $5\,v_\mathrm{th}$).
A box that is too small is reported with the $|v|_{\max}$ to re-run with.


\subsection{Multispecies Outputs (Version 4.0)}
\label{sec:out_ms}

\paragraph{Per-species files.}  With \texttt{i\_multi = -1}, the files of a species carry a tag after
the case name: none for the heated species, \texttt{\_e} for the electrons, \texttt{\_i1},
\texttt{\_i2}, \dots\ for the bulk ions (e.g.\ \texttt{fout-N1\_e.txt}).
\begin{description}
\item[\texttt{finit-*.txt}] Initial distribution of every species on its grid
  ($\vp$, $\vpa$, $f$).
\item[\texttt{fout-*.txt}, \texttt{xout-*.dat}] Final distribution of every species (coupled run),
  same layout as in v3.2.  The grid of a regridded species is the final one.
\item[\texttt{power\_lossE\_vs\_time.txt}] (v3.2 path, $\tau_E>0$) time, $P_E$ (MW/m$^3$) and
  $-P_E\tau_E/W$ of the heated species.
\end{description}

\paragraph{Coupled-loop traces} (one line per step; species in the order heated, electrons, bulk ions;
pairs $(s\leftarrow j)$ with $s$ outer and $j$ inner):
\begin{description}
\item[\texttt{ms\_moments\_vs\_time.txt}] Per species: $n$ (m$^{-3}$), $W$ (J/m$^3$),
  $T_\mathrm{eff}$, $T_\perp$, $T_\parallel$ (eV), $P_\parallel$ (kg\,m$^{-2}$\,s$^{-1}$); then $W_\mathrm{tot}$ and
  $P_{\parallel,\mathrm{tot}}$.  The first line is the initial state.  For a drifting species the
  thermal temperature is $T_\mathrm{eff}-mu^2/3$, $u=P_\parallel/nm$.
\item[\texttt{ms\_power\_pairs\_vs\_time.txt}] Raw power of every pair (MW/m$^3$), before the
  projection, and their sum.
\item[\texttt{ms\_momentum\_pairs\_vs\_time.txt}] Raw momentum rate of every pair
  (kg\,m$^{-2}$\,s$^{-2}$), and their sum.
\item[\texttt{ms\_density\_pairs\_vs\_time.txt}] Raw particle rate of every pair
  (m$^{-3}$\,s$^{-1}$).
\item[\texttt{ms\_power\_nrl\_vs\_time.txt}] Maxwellian (NRL) exchange power of every pair at the same
  temperatures and Coulomb logarithms.
\item[\texttt{ms\_coulomb\_log\_vs\_time.txt}] Coulomb logarithm of every pair.
\item[\texttt{ms\_power\_other\_vs\_time.txt}] $P_\mathrm{RF}$ and $P_{E,s}$ of every species (MW/m$^3$),
  on the corrected state.
\item[\texttt{ms\_source\_vs\_time.txt}] Beam source and particle loss: $\mathrm dN/\mathrm dt$,
  power, $\mathrm dP_\parallel/\mathrm dt$ of each; then the momentum rate of $L_E$ for every species.
\item[\texttt{ms\_corrections\_vs\_time.txt}] Per species: the defects removed by the projection
  ($\mathrm dN$, $\mathrm dP$ in MW/m$^3$, $\mathrm dR$) and the three projection amplitudes.
\item[\texttt{ms\_regrid.txt}] One line per regrid: time, species, $T$, $u$, old and new half-width and
  centre, the moment errors of the transfer before the restoration, the restoration amplitudes.
\end{description}

\paragraph{Standard output.}  The species and grids table at start; the projection and Picard
settings; grid warnings and regrid events; every 50 steps the last Picard change; at the end a
wall-clock profile of the coupled step, the final state of every species and the PARDISO statistics
of every species.

\paragraph{Checking a coupled run.}  From \texttt{ms\_moments}: the particle numbers and, without
source, $W_\mathrm{tot}$ and $P_{\parallel,\mathrm{tot}}$ must be constant to $\sim10^{-10}$.  With a source,
$(W^{n+1}_\mathrm{tot}-W^n_\mathrm{tot})/\Delta t$ must equal the sum of $P_\mathrm{RF}$, the source and loss
powers and $\sum_sP_{E,s}$ of the same step, and at steady state $P_\mathrm{RF}\tau_E/\sum_sW_s\to1$
(RF) or $P_\mathrm{NBI}=\sum_sW_s/\tau_E+W_a/\tau_s$ (NBI).


% ============================================================
\section{Example Namelist}
% ============================================================

The following example illustrates a time-dependent D-beam simulation in a D
plasma with nonlinear self-collisions, using a two-phase time-stepping
strategy: 1000 coarse steps followed by 1000 finer steps. It lists \emph{every}
\texttt{\&INPUT} parameter; the RF block is included for completeness but is only
read when \texttt{irf=-1}, and the diagnostic/advanced flags are shown at their
default (off) values.

A second, fully commented example ships with the code as
\texttt{example\_namelist\_ITER\_N2.txt} (project root): the ITER $N=2$ RF
reference case on $150\times150$ with nonlinear self-collisions, which converges
in 900 steps (about 21~s on 16 threads).  It is the quickest way to check a new
build:
\begin{lstlisting}
FP2D_QLRF_NL.exe < example_namelist_ITER_N2.txt
\end{lstlisting}

\begin{lstlisting}
&INPUT
  casename = 'JET-beam7-TD0-NLSC',

  ! Grid
  nperp = 151, npar = 151,
  vperp_min = 0.d0,  vperp_max = 5.d6,
  vpar_min  = -4.d6, vpar_max  = 4.d6,
  ising = 1, nsing = 25, vbound = 0.07d6,
  p_grid = 2.d0,

  ! Plasma: electrons + D ions
  nbulk = 2,
  t     = 15.d3, 10.d3, 8*0.d0,
  aa = 2.d0, za = 1.d0,
  ab = 2.d0, 8*0.d0,
  zb = 1.d0, 8*0.d0,
  ne     = 5.0d19,
  xpart  = 0.05d0,
  xb     = 0.50d0, 8*0.d0,

  ! D-NBI beam at 120 keV, 90 degrees
  isource = -1,
  beam_ekin      = 120.d0,
  beam_angle_deg = 90.d0,
  beam_dvperp    = 0.12d6,
  beam_dvpar     = 0.20d6,
  taus = 300.d-3,

  ! RF quasi-linear operator (b0..delta_RF read only when irf = -1)
  irf   = 0,
  b0    = 3.45d0,
  frek  = 53.d6,
  nharm = 1,
  kpar  = 2.d0,
  kperp = (30.d0, 1.d0),
  eplus = (1.d3, 3.25d3),
  emin  = (3.d3, 3.75d3),
  delta_RF = 0.5d6,

  ! Time integration: two phases
  icn      = 1,
  ntimes   = 1000, 1000, 0,
  timestep = 1.d-3, 10.d-3, 0.d0,
  iold   = 0,
  istart = 0,
  isc    = -1,
  core_frac = 3.8d-3,

  ! Steady-state convergence detection
  i_ss_check     = -1,
  n_ss_window    = 50,
  ss_tol_eps     = 1.d-3,
  ss_tol_tail    = 1.d-2,
  ss_tol_moment  = 1.d-3,
  i_conv_shape   = 0,

  ! Output
  iplot_pow = -1,
  iplot_mom = -1,
  notxt     = 0,
  n_snap    = 0,

  ! Advanced numerics (defaults shown)
  i_upwind     = 0,
  i_lu_reuse   = -1,
  lu_reuse_tol = 1.d-10,
  n_threads    = 0,      ! 0: 16, or MKL_NUM_THREADS; n > 0: n threads
/
\end{lstlisting}


\subsection{Multispecies Examples (Version 4.0)}

\paragraph{RF heating with evolving backgrounds.}  ITER fundamental tritium heating in a 50:50 D--T
plasma ($N=1$ benchmark), all species at 12\,keV initially, $\tau_E=0.5$\,s on every species; steady
state within a few seconds at $T_\mathrm T=22.3$, $T_\mathrm D=17.9$, $T_e=11.4$\,keV.  Only the entries that
differ from the v3.2 example are shown.
\begin{lstlisting}
  casename = 'N1',
  nperp = 150, npar = 150,
  vperp_min = 0.d0, vperp_max = 2.d7, vpar_min = -6.d6, vpar_max = 6.d6,
  nbulk = 2, t = 12.d3, 12.d3, 8*0.d0,
  aa = 3.d0, za = 1.d0, ab = 2.d0, 8*0.d0, zb = 1.d0, 8*0.d0,
  ne = 15.d19, xpart = 50.d-2, xb = .50d0, 8*0.d0,
  isource = 0,
  irf = -1, eplus = (2.5d2,1.0d2), emin = (7.5d2,1.0d0), kperp = (30.d0,0.d0),
  kpar = 2.d0, frek = 30.45d6, delta_RF = 2.5d6, b0 = 6.d0, nharm = 1,
  icn = 1, ntimes = 100, 300, 300, timestep = 5.d-3, 10.d-3, 25.d-3,
  iold = 0, istart = 1, isc = -1, n_threads = 8,
  ! version 4.0
  i_multi = -1, i_coupled = -1, ta_init = 12.d3,
  lmax_leg = 16, i_ms_cons = 2, n_picard = 1,
  vbox_b = 8.d0, 12.d0, tau_E = 0.5d0
\end{lstlisting}

\paragraph{Neutral-beam injection.}  JET beam case 1: deuterium beam at 120\,keV and $45^\circ$,
$\tau_s=0.2$\,s, beam density $0.15\,n_e$, into deuterons at 10\,keV and electrons at 15\,keV,
$\tau_E=1$\,s on every species, starting from an empty beam distribution.  With regridding, the bulk
grids follow the cooling electrons and the drifting deuterons.
\begin{lstlisting}
  casename = 'B1',
  nperp = 151, npar = 151,
  vperp_min = 0.d0, vperp_max = 6.d6, vpar_min = -5.d6, vpar_max = 6.d6,
  nbulk = 2, t = 15.d3, 10.d3, 8*0.d0,
  aa = 2.d0, za = 1.d0, ab = 2.d0, 8*0.d0, zb = 1.d0, 8*0.d0,
  ne = 5.d19, xpart = 15.d-2, xb = .85d0, 8*0.d0,
  isource = -1, beam_ekin = 120.d0, beam_angle_deg = 45.d0,
  beam_dvperp = 0.12d6, beam_dvpar = 0.2d6, taus = 200.d-3,
  irf = 0,
  icn = 1, ntimes = 200, 600, 0, timestep = 5.d-3, 25.d-3, 25.d-3,
  iold = 0, istart = 0, isc = -1, n_threads = 8,
  ! version 4.0
  i_multi = -1, i_coupled = -1, n_picard = 1, lmax_leg = 16, i_ms_cons = 2,
  vbox_b = 8.d0, 12.d0, tau_E = 1.d0, i_regrid = -1
\end{lstlisting}

\paragraph{Equilibration without source.}  Without RF or beam (\texttt{irf = 0},
\texttt{isource = 0}) the coupled loop relaxes all species to a common temperature at constant total
energy; $\tau_E$ is ignored.  With $T_e=20$\,keV, $T_\mathrm T=T_\mathrm D=10$\,keV
(\texttt{t = 20.d3, 10.d3}, \texttt{ta\_init = 10.d3}) the temperatures follow the Maxwellian
two-temperature model to 0.1--0.5\% of their change and settle at 15.0\,keV within about 2\,s.


% ============================================================
\section{Compilation and Execution}
% ============================================================

\subsection{Build}

Open the Visual Studio solution
\texttt{FP2D\_QLRF\_NL/FP2D\_QLRF\_NL.sln} and build the
\texttt{Release|x64} configuration.  The compiler is Intel \texttt{ifx}
with Intel MKL; the Release build links the \emph{threaded} MKL, so the
sparse solver PARDISO runs in parallel, and is compiled with OpenMP, which
parallelises the code's own loops over the grid: the matrix assembly, the
self-collision potentials and the per-step diagnostics (the Debug build is
sequential).  The executable is placed at
\texttt{FP2D\_QLRF\_NL/x64/Release/FP2D\_QLRF\_NL.exe}.

\paragraph{Version 4.0 sources.}  The main program \texttt{FP\_Coll\_2D\_ms} is in
\texttt{main-FP\_Coll\_2D.f90}.  The multispecies extension adds four files to the project:
\texttt{species.f90} (species containers and grids), \texttt{legendre\_pot.f90} (Legendre
potentials and pair coefficients), \texttt{ms\_checks.f90} (checks, \texttt{i\_ms\_check}) and
\texttt{TimeFP\_ms.f90} (coupled time loop, projection and regridding).  The method of every routine is
described stage by stage in the appendix of the multispecies article.

\paragraph{Number of threads.}\label{sec:threads}  The code uses 16 threads by default (fewer on a
machine with fewer cores) and prints the number at start.  To change it, set
\texttt{n\_threads} in the namelist: \texttt{n\_threads = 4} uses four threads
for both MKL (PARDISO) and the OpenMP loops, and overrides the environment.
With \texttt{n\_threads = 0} (the default) the environment variable
\texttt{MKL\_NUM\_THREADS} (or \texttt{OMP\_NUM\_THREADS}) is used if set before
running, e.g.\ \texttt{set MKL\_NUM\_THREADS=4} in a Windows command prompt, and
16 otherwise.  There is no separate sequential solver: \texttt{n\_threads = 1}
is the sequential code, with results identical to the sequential (Debug) build
and reproducible from run to run.  The thread count is chosen at run time; only
the choice between threaded and sequential MKL is fixed at compilation (Release
and Debug respectively).  MKL's own default, one thread per physical core, is slower than 8--16
threads on a many-core machine, which is why the code sets its own.  The OpenMP
loops use the same number unless \texttt{OMP\_NUM\_THREADS} is set.  On a
$150\times150$ grid with nonlinear self-collisions and RF, 100 steps took 29.4~s
with the original sequential code and take 3.0~s now on 16 threads (with LU
reuse, below); a full 900-step run takes 21~s.  The gain saturates at 8--16
threads.  The OpenMP loops do not change the results; with one thread the
results are identical to the sequential build, and with more the parallel
factorization in MKL makes them differ in the last digits (up to
${\sim}10^{-5}$ of $\max f$ after 100 steps), because it eliminates in a
different order.  When running several cases at once, give each only a few
threads.

\paragraph{LU reuse.}\label{sec:lureuse}  The matrix $I - \theta\,\Delta t\,L$
changes from step to step, but slowly.  With \texttt{i\_lu\_reuse = -1} (the
default) the code therefore factorises it only occasionally: at the other steps
it solves by a few CGS iterations preconditioned with the last LU factors
(PARDISO phase~23), checks the true residual of the result, and refactorises
only if that residual exceeds
$\max(\texttt{lu\_reuse\_tol},\,10\,r_\mathrm{direct})$.  The end of the run
prints how many factorisations were needed and the largest residuals, e.g.
\begin{lstlisting}
PARDISO: 900 steps, 2 factorisations; LU reused on 898 steps
         (2.8 CGS iterations on average, 1 rejected)
\end{lstlisting}
On the ITER $N=2$ reference case this cut a full run from 142~s to 52~s and on a
JET beam fill-up from 96~s to 28~s (16 threads); with one thread, 100 steps went
from 28.7~s to 8.3~s.  The results agree with refactorising every step to
${\sim}10^{-8}$ of $\max f$ on one thread, and within the threading noise on
16.  The finite-difference weights, which depend only on the grid, are also
computed once per run instead of at every step.

\subsection{Running}

The program reads from \texttt{stdin}:
\begin{lstlisting}
FP2D_QLRF_NL.exe < inputs/jet_beam_case1.dat
\end{lstlisting}
Output files are written to the working directory.

\subsection{Plotting}
\label{sec:plotting}

Use the companion Python script \texttt{fp2d\_plot.py}:
\begin{lstlisting}
python fp2d_plot.py run FP2D_QLRF_NL.exe inputs/jet_beam_case1.dat --show
python fp2d_plot.py plot x64/Release --casename JET-beam7-TD0-NLSC --save plots/
python fp2d_plot.py compare x64/Release --cases JET-beam7-TD0-Lin JET-beam7-TD0-NLSC
python fp2d_plot.py plot x64/Release --cases JET-beam7-TD0-NLSC --movie
\end{lstlisting}
The last line builds an animation of $f$, the kinetic-energy density and its
perpendicular and parallel parts (a $2\times2$ array of panels) from the \texttt{vdf\_snap} files (requires \texttt{n\_snap > 0}); \texttt{--movie-scale
frame} rescales each frame to its own maximum, and \texttt{--movie-format}
selects AVI (default), MP4 or GIF.

The axes of \texttt{plot} and \texttt{compare} figures can be set with
\texttt{--logx}, \texttt{--logy} (logarithmic axes) and
\texttt{--xrange xmin:xmax}, \texttt{--yrange ymin:ymax} (zoom).
\texttt{--logf} puts the plotted quantity on a log scale: the colour scale of a
2D map, and the $y$-axis of a 1D profile or of a time trace.  An axis on which
the data really change sign ($\vpa$, or a power) is left linear, with a note.
See \texttt{fp2d\_plot-readme.txt} for the full documentation of the
plotting tool.

\end{document}
